% ============================================================
% Corpo das questões — GERADO por gerar-corpo.py a partir de
% conteudo-blocos.json (SSOT). Não editar à mão: editar o JSON
% e regenerar. Sincronia prova/solução garantida via \ifsolucao
% ============================================================

\ifsolucao\else
\begin{center}
\renewcommand{\arraystretch}{1.3}
\begin{tabular}{|l|l|c|c|}
\hline
\textbf{Bloco} & \textbf{Conteúdo} & \textbf{Aulas} & \textbf{Pontos} \\ \hline
1 & Teoria do Consumidor & 1--3 & 25 \\ \hline
2 & Equilíbrio Geral & 4--6 & 25 \\ \hline
3 & Externalidades e Bens Públicos & 7 & 25 \\ \hline
4 & Assimetria Informacional e Sinalização & 8--9 & 25 \\ \hline
\multicolumn{3}{|r|}{\textbf{Total}} & \textbf{100} \\ \hline
\end{tabular}
\end{center}
\vspace{0.5em}
\fi


\section*{Bloco 1 --- Teoria do Consumidor \hfill \normalsize (25 pontos)}

\textbf{Questão 1.} (10 pontos) \textbf{Itens objetivos.} Nos itens \textbf{a)} e \textbf{b)}, assinale a \emph{única} alternativa correta (3 pontos cada; não há penalização por erro). Nos itens \textbf{c)} e \textbf{d)}, responda \textbf{Verdadeiro ou Falso} (2 pontos cada; vale a regra de anulação descrita na capa). Registre suas respostas no quadro-resposta ao final do bloco.

\textbf{a)} (3 pontos) Em maio de 2026, o governo brasileiro extinguiu a alíquota federal de $20\%$ do imposto de importação sobre remessas de até US\$~50 --- a chamada ``taxa das blusinhas'' ---, barateando as compras em plataformas internacionais. Considere um consumidor para quem as peças de vestuário importadas de baixo custo ($x_1$) constituem um \emph{bem inferior} (ao enriquecer, migra para marcas nacionais). Num dado ponto, estimam-se o efeito-preço total marshalliano $\dfrac{\partial x_1^{M}}{\partial p_1}=-3$, a quantidade $x_1^{M}=20$ e a propensão marginal a consumir $\dfrac{\partial x_1^{M}}{\partial m}=-0{,}2$. Aplicando a equação de Slutsky para o próprio preço, $\dfrac{\partial h_1}{\partial p_1}=\dfrac{\partial x_1^{M}}{\partial p_1}+x_1^{M}\dfrac{\partial x_1^{M}}{\partial m}$, qual é o valor do efeito-substituição $\dfrac{\partial h_1}{\partial p_1}$ nesse ponto?

\begin{enumerate}[label=\Alph*., leftmargin=3.2em, itemsep=0.15em, topsep=0.2em]
  \item $-3$
  \item $-7$
  \item $-4$
  \item $+1$
  \item $+4$
\end{enumerate}
\gabaritoitem{B}
\begin{solucao}
{\footnotesize\textit{Paralelo de: Simulado 3, Bloco 1, ME2 (Azul-Gol / Slutsky para bem inferior).}}\par\vspace{0.2em}
\textbf{Passo 1 --- o que se pede.} A equação de Slutsky decompõe o efeito-preço total marshalliano em efeito-substituição (hicksiano) mais efeito-renda. Como o bem é inferior ($\partial x_1^{M}/\partial m<0$), o efeito-renda tem sinal \emph{oposto} ao do caso normal, e o efeito-substituição não coincide com o efeito total. Pede-se isolar $\partial h_1/\partial p_1$.

\textbf{Passo 2 --- montar Slutsky.} Substituindo os valores dados,
\[
\frac{\partial h_1}{\partial p_1}=\frac{\partial x_1^{M}}{\partial p_1}+x_1^{M}\frac{\partial x_1^{M}}{\partial m}=-3+20\cdot(-0{,}2)=-3-4=-7.
\]

\textbf{Passo 3 --- leitura das três peças.} O efeito-substituição $(-7)$ é \emph{mais negativo} que o efeito total $(-3)$. A razão: o efeito-renda \emph{sobre a demanda} é $-x_1^{M}\,\partial x_1^{M}/\partial m=-20\cdot(-0{,}2)=+4$, isto é, positivo. Conferindo a soma: total $(-3)=$ substituição $(-7)+$ renda $(+4)$. O sinal negativo de $\partial h_1/\partial p_1$ respeita a lei da demanda compensada, como deve.

\textbf{Passo 4 --- distratores tentadores.} O distrator $+1$ é o mais sedutor: vem de tratar o bem como \emph{normal}, trocando o sinal da propensão marginal ($-3+20\cdot 0{,}2=+1$); quem o marca esqueceu que a inferioridade ($\partial x_1^{M}/\partial m<0$) faz o termo de Slutsky \emph{somar} negatividade em vez de subtraí-la. O distrator $-3$ é o segundo mais tentador: confunde o efeito-preço total marshalliano com o efeito-substituição, ignorando que num bem inferior eles \emph{não coincidem} --- erro de quem aplica a intuição do bem normal, em que substituição e renda puxam no mesmo sentido. O distrator $-4$ reporta isoladamente o termo $x_1^{M}\,\partial x_1^{M}/\partial m=-4$, sem somar o efeito total; o distrator $+4$ é o efeito-renda \emph{sobre a demanda} ($+4$), de sinal correto mas que não é o que se pede.

\textbf{Intuição econômica:} encarecida a peça importada, o efeito-substituição puro a torna relativamente menos atraente e empurra a quantidade para baixo. Mas, sendo um bem inferior, o encarecimento também empobrece o consumidor em termos reais, e esse empobrecimento o faz querer \emph{mais} do bem barato (efeito-renda positivo sobre a quantidade), o que \emph{atenua} a queda observada. Logo a substituição pura tem de ser ainda mais forte ($-7$) que o efeito total ($-3$). O fim da ``taxa das blusinhas'' baixou o preço dessas peças: para um bem inferior, parte de quem voltou a comprá-las o fez não por gosto, mas porque a economia no orçamento funcionou como um leve enriquecimento --- exatamente o canal de renda que aqui mascara a força da substituição.
\end{solucao}

\textbf{b)} (3 pontos) A Smart Fit fechou 2025 com mais de 2.000 academias e receita superior a R\$~7 bilhões, consolidando a liderança do fitness de baixo custo na América Latina. Um analista modela a demanda de um cliente a partir da \emph{utilidade indireta} estimada $v(p_1,p_2,m)=\dfrac{m}{p_1^{3/4}\,p_2^{1/4}}$, em que $x_1$ é o número de meses de acesso à academia contratados e $x_2$ um bem composto (preços $p_1,p_2$ e renda $m$ em unidades monetárias padronizadas). O expoente $3/4$ sobre $p_1$ reflete um perfil de cliente para quem a atividade física absorve fração elevada do orçamento de bem-estar. Aplicando a \emph{identidade de Roy} para recuperar a demanda marshalliana por academia e avaliando-a em $m=120$ e $p_1=3$, qual é a quantidade demandada de meses de academia?

\begin{enumerate}[label=\Alph*., leftmargin=3.2em, itemsep=0.15em, topsep=0.2em]
  \item $10$
  \item $40$
  \item $20$
  \item $60$
  \item $30$
\end{enumerate}
\gabaritoitem{E}
\begin{solucao}
{\footnotesize\textit{Paralelo de: Simulado 2, Bloco 1, ME2 (Três Corações / identidade de Roy a partir de utilidade indireta).}}\par\vspace{0.2em}
\textbf{Passo 1 --- o que se pede.} A identidade de Roy recupera a demanda marshalliana a partir da utilidade indireta sem reabrir a UMP:
\[
x_1^{M}(p,m)=-\frac{\partial v/\partial p_1}{\partial v/\partial m}.
\]

\textbf{Passo 2 --- derivadas parciais.} De $v=m\,p_1^{-3/4}p_2^{-1/4}$,
\[
\frac{\partial v}{\partial p_1}=m\left(-\tfrac34\right)p_1^{-7/4}p_2^{-1/4},\qquad \frac{\partial v}{\partial m}=p_1^{-3/4}p_2^{-1/4}.
\]

\textbf{Passo 3 --- aplicar Roy.} A razão cancela $p_2^{-1/4}$ e parte de $p_1$:
\[
x_1^{M}=-\frac{m\left(-\tfrac34\right)p_1^{-7/4}p_2^{-1/4}}{p_1^{-3/4}p_2^{-1/4}}=\frac34\,m\,p_1^{-7/4+3/4}=\frac34\cdot\frac{m}{p_1}=\frac{3m}{4p_1}.
\]

\textbf{Passo 4 --- avaliação e distratores.} Em $m=120,\ p_1=3$: $x_1^{M}=\dfrac{3\cdot 120}{4\cdot 3}=\dfrac{360}{12}=30$. O distrator $10$ é o erro de quem usa o expoente $1/4$ de $p_1$ como participação orçamentária (em vez de $3/4$), obtendo $\tfrac14\cdot\tfrac{m}{p_1}=10$: é a troca clássica entre o expoente do bem e o do bem composto. O distrator $40$ resulta de esquecer o fator de participação e reportar $m/p_1=40$ --- aplica-se Roy mecanicamente, mas perde-se o coeficiente $\tfrac34$ no cancelamento dos expoentes. O distrator $20$ assume simetria Cobb-Douglas $\tfrac12\cdot\tfrac{m}{p_1}=20$, e o $60$ esquece o preço no denominador ($m/2$).

\textbf{Intuição econômica:} a utilidade indireta já contém todo o comportamento ótimo --- a identidade de Roy apenas o ``desempacota'', derivando-a em preço e renda. O expoente $3/4$ sobre $p_1$ sinaliza que a academia absorve três quartos do orçamento de bem-estar desse cliente, o que casa com o perfil fitness intensivo que sustenta a expansão da Smart Fit; multiplicado por $m/p_1$, isso dá a quantidade física (meses) contratada. É a face dual do problema: em vez de maximizar utilidade sujeito ao orçamento, lê-se a demanda diretamente da função-valor.
\end{solucao}

\textbf{c)} (2 pontos, V ou F) Com a fusão entre Petz e Cobasi aprovada pelo CADE em 2025, o setor pet --- que movimenta dezenas de bilhões de reais ao ano no Brasil --- voltou ao centro do debate. Considere um tutor que aloca todo o orçamento destinado ao seu animal entre vários itens (ração, medicamentos, banho \& tosa, hospedagem em hotelzinho, entre outros) e suponha que os \emph{serviços premium} (banho \& tosa e hospedagem) sejam um \emph{bem de luxo}, isto é, tenham elasticidade-renda maior que $1$ (quando a renda sobe $1\%$, o gasto com esses serviços sobe mais de $1\%$). Sob as convenções do curso (preferências localmente não saciadas, orçamento integralmente exaurido), pode-se afirmar que \emph{necessariamente} pelo menos um outro item da cesta pet tem elasticidade-renda menor que $1$.

\gabaritoitem{V}
\begin{solucao}
{\footnotesize\textit{Paralelo de: Simulado 2, Bloco 1, VF1 (carne bovina / agregação de Engel: luxo implica necessariamente que algum outro bem tem elasticidade-renda < 1).}}\par\vspace{0.2em}
\textbf{Passo 1 --- a identidade de agregação de Engel.} Diferenciando a restrição orçamentária $\sum_i p_i x_i^{M}(p,m)=m$ em relação a $m$ obtém-se $\sum_i p_i \dfrac{\partial x_i^{M}}{\partial m}=1$. Multiplicando e dividindo cada termo por $x_i^{M}$ e por $m$, isso vira a \emph{agregação de Engel}:
\[
\sum_i s_i\,\eta_i=1,\qquad s_i=\frac{p_i x_i^{M}}{m}\ \ (\text{participações, com }\textstyle\sum_i s_i=1),\quad \eta_i=\frac{\partial x_i^{M}}{\partial m}\frac{m}{x_i^{M}}.
\]

\textbf{Passo 2 --- o argumento.} A média das elasticidades-renda, \emph{ponderada pelas participações}, é exatamente $1$. Se os serviços premium têm $\eta>1$, eles puxam a média ponderada para cima. Para que a média volte a $1$, é impossível que \emph{todos} os demais itens tenham $\eta\ge 1$: nesse caso $\sum_i s_i\eta_i>1$, contradição (com participações positivas). Logo existe ao menos um item com $\eta_i<1$.

\textbf{Passo 3 --- por que vale ``necessariamente''.} A conclusão não depende da forma funcional nem dos números: é consequência puramente contábil de o orçamento ser exaurido. Basta haver mais de um item com participação positiva. Note que o teorema garante a \emph{existência} de algum item com $\eta_i<1$, sem identificá-lo: não se afirma que seja a ração, o medicamento ou qualquer item específico. Por isso a afirmação é \textbf{verdadeira}.

\textbf{Intuição econômica:} renda extra precisa ir para algum lugar, e a soma dos gastos sempre fecha em $100\%$. Se o tutor destina uma fatia \emph{mais que proporcional} do aumento de renda a serviços premium (luxo), sobra menos que proporcional para o conjunto dos outros itens --- algum deles cresce abaixo da renda (necessidade) ou até encolhe (inferior). A ração básica é a candidata natural: quando o tutor enriquece, mantém a alimentação essencial mas gasta proporcionalmente mais em mimos. Não pode haver cesta ``só de luxos''; luxos e necessidades coexistem por força da restrição orçamentária.
\end{solucao}

\textbf{d)} (2 pontos, V ou F) Com a multiplicação dos planos de streaming com anúncios (Disney+ e Max passaram a oferecer camadas mais baratas com publicidade), um assinante ordena planos descritos por pares $(q,a)$, em que $q$ é a qualidade máxima de vídeo e $a$ a ausência de anúncios (quanto maior $a$, menos interrupções). Ele os compara \emph{lexicograficamente}: prefere estritamente o plano de maior $q$ e, \emph{somente} em caso de empate em $q$, usa $a$ como desempate (maior $a$ preferido). Formalmente, $(q,a)\succeq(q',a')$ se $q>q'$, ou se $q=q'$ e $a\ge a'$. Como essa relação de preferências é completa e transitiva, pode-se afirmar que ela é \emph{necessariamente} representável por uma função utilidade contínua $U:\mathbb{R}^2\to\mathbb{R}$.

\gabaritoitem{F}
\begin{solucao}
{\footnotesize\textit{Paralelo de: Simulado 1, Bloco 1, VF2 (B3 / ordem lexicográfica: completa e transitiva mas NÃO representável por utilidade contínua --- argumento de Debreu).}}\par\vspace{0.2em}
\textbf{Passo 1 --- o que se afirma.} A sentença pretende que completude $+$ transitividade bastem para garantir representação por utilidade contínua. Vamos exibir um contraexemplo: a própria preferência lexicográfica do assinante é completa e transitiva, mas \emph{não admite representação por utilidade alguma} --- muito menos contínua.

\textbf{Passo 2 --- completa e transitiva, sim.} Dados dois planos quaisquer, ou suas qualidades diferem (decide $q$), ou empatam (decide $a$): todo par é comparável, logo a relação é completa. A transitividade segue de encadear as comparações na ordem $q$-primeiro, $a$-depois. Portanto as hipóteses do enunciado valem.

\textbf{Passo 3 --- não há utilidade que a represente.} Suponha, por absurdo, que exista $U$ com $(q,a)\succeq(q',a')\iff U(q,a)\ge U(q',a')$. Fixe $q$ e compare $(q,1)\succ(q,0)$: então $U(q,1)>U(q,0)$. Para cada $q$ escolha um racional $r(q)\in\big(U(q,0),\,U(q,1)\big)$. Se $q'>q$, todo plano com qualidade $q'$ é estritamente preferido a todo plano com qualidade $q$, logo $U(q',0)>U(q,1)$, o que força $r(q')>r(q)$. Assim $q\mapsto r(q)$ é uma injeção \emph{estritamente crescente} de $\mathbb{R}$ nos racionais $\mathbb{Q}$ --- impossível, pois $\mathbb{R}$ é não-enumerável e $\mathbb{Q}$ é enumerável. Contradição: nenhuma $U$ existe. (É o argumento clássico de Debreu para a não-representabilidade da ordem lexicográfica.)

\textbf{Passo 4 --- o que realmente faltou.} A representabilidade por utilidade \emph{contínua} exige, além de completude e transitividade, a \emph{continuidade} das preferências (conjuntos $\succeq$ e $\preceq$ fechados). A lexicográfica viola exatamente essa continuidade: o conjunto dos planos preferidos a um ponto não é fechado. Logo a afirmação é \textbf{falsa}.

\textbf{Intuição econômica:} um assinante que coloca a qualidade de vídeo em primeiro lugar de forma \emph{absoluta} --- nenhuma dose de ``sem anúncios'' compensa um pingo a menos de qualidade --- não admite ``taxa de troca'' entre os dois atributos. Mas uma função utilidade é justamente o que codifica trocas marginais: ela atribui um número que permite compensar um atributo pelo outro. Quando a prioridade é lexicográfica, não existe número que respeite essa ordem rígida sobre um contínuo de qualidades. Completude e transitividade garantem uma \emph{ordenação} coerente, não que ela caiba dentro da reta real de forma contínua.
\end{solucao}

\quadroresposta{1}

\textbf{Questão 2.} (15 pontos) A alta histórica do cacau em 2024--2025 --- com a tonelada superando US\$~12 mil em Nova York --- pressionou os custos da indústria de chocolate; a A alta pressionou as margens da indústria nacional de chocolate. Suponha que, diante desse cenário, o governo cogite \emph{hipoteticamente} um imposto \emph{específico} sobre o chocolate fino importado. Considere um consumidor que reparte sua renda mensal entre chocolate fino importado ($x_1$, em barras) e um bem composto ($x_2$, numerário, com $p_2=1$), segundo preferências $u(x_1,x_2)=x_1^{1/2}\,x_2^{1/2}$. Sua renda disponível para essas categorias é $m=R\$~72$ e o preço da barra, \emph{sem} imposto, é $p_1^{0}=R\$~1$. Suponha que o imposto específico seja $t=R\$~8$ por barra, repassado integralmente ao consumidor, de modo que o preço passe a $p_1^{1}=p_1^{0}+t=R\$~9$. A receita arrecadada é $R_{\text{arr}}=t\,x_1$, avaliada na quantidade efetivamente consumida sob o imposto. Em todos os itens, considere solução interior.

\textbf{a)} (5 pontos) Monte o problema de maximização da utilidade (UMP) e derive as demandas marshallianas $x_1^{M}(p,m)$ e $x_2^{M}(p,m)$ em forma fechada para esta Cobb-Douglas simétrica. Em seguida, calcule as cestas consumidas \emph{antes} do imposto ($p_1^{0}=1$) e \emph{depois} do imposto ($p_1^{1}=9$), e obtenha a receita $R_{\text{arr}}=t\,x_1$ efetivamente arrecadada.
\resolucaobox{7.2cm}
\begin{solucao}
{\footnotesize\textit{Paralelo de: Simulado 3, Bloco 1, Aberta (suco de laranja / Cobb-Douglas: tarifa específica, CV, lump-sum e peso-morto).}}\par\vspace{0.2em}
\textbf{Passo 1 --- UMP e CPO.} O consumidor resolve $\max\,x_1^{1/2}x_2^{1/2}$ s.a. $p_1x_1+x_2=m$ (com $p_2=1$). Igualando a $\text{TMS}$ à razão de preços, $\dfrac{x_2}{x_1}=\dfrac{p_1}{p_2}=p_1$, ou seja $x_2=p_1x_1$.

\textbf{Passo 2 --- forma fechada.} Levando ao orçamento $p_1x_1+p_1x_1=m$, obtêm-se as demandas Cobb-Douglas com participação $1/2$ em cada bem:
\[
x_1^{M}(p_1,m)=\frac{m}{2p_1},\qquad x_2^{M}(p_1,m)=\frac{m}{2}.
\]

\textbf{Passo 3 --- cestas antes e depois.} Com $m=72$:
\[
\text{antes }(p_1^{0}=1):\quad x_1=\frac{72}{2}=36,\quad x_2=\frac{72}{2}=36;
\]
\[
\text{depois }(p_1^{1}=9):\quad x_1=\frac{72}{18}=4,\quad x_2=\frac{72}{2}=36.
\]
Ambas interiores. Note que $x_2$ não muda: na Cobb-Douglas, o gasto em cada bem é fração fixa da renda (R\$~36 em cada), então a alta de $p_1$ derruba a \emph{quantidade} de chocolate, mas preserva o gasto.

\textbf{Passo 4 --- receita.} O imposto incide sobre o chocolate efetivamente consumido sob $p_1^{1}=9$, isto é, $x_1=4$:
\[
R_{\text{arr}}=t\,x_1=8\cdot 4=32.
\]

\textbf{Intuição econômica:} o imposto de R\$~8 por barra multiplica por nove o preço de varejo e corta o consumo de $36$ para $4$ barras. Como a Cobb-Douglas mantém constante a fatia do orçamento gasta em cada bem, o consumidor continua destinando R\$~36 ao chocolate --- só que agora compra muito menos volume por esse mesmo valor. A arrecadação efetiva, R\$~32, já embute essa retração da quantidade: é bem menor do que seria se o consumo tivesse ficado parado em $36$ barras.
\rubricalabel
\begin{itemize}
\item Monta a UMP corretamente (objetivo $+$ restrição com $p_2=1$): 1 pt.
\item Iguala $\text{TMS}=p_1/p_2$ e obtém a forma fechada $x_1^{M}=m/(2p_1)$, $x_2^{M}=m/2$: 1{,}5 pt.
\item Calcula a cesta antes $(36,36)$ e depois $(4,36)$: 1 pt (0{,}5 cada).
\item Obtém a receita $R_{\text{arr}}=t\,x_1=8\cdot 4=32$, avaliada na quantidade pós-imposto (não na pré): 1{,}5 pt.
\item Avaliar $R_{\text{arr}}$ na quantidade \emph{pré}-imposto ($8\cdot 36=288$): perde os 1{,}5 pt da receita, mantém o restante.
\item Erro aritmético isolado com método integralmente correto: penalidade leve de $-0{,}5$ pt (não zera o item).
\end{itemize}
\end{solucao}

\textbf{b)} (5 pontos) Obtenha a função utilidade indireta $v(p_1,m)$ desta Cobb-Douglas e, invertendo-a em $m$, a função dispêndio $e(p_1,\bar u)$ (lembre que, pelo lema de Shephard, $\partial e/\partial p_1=h_1$). Calcule os níveis de utilidade $\bar u^{0}=v(p_1^{0},m)$ e $\bar u^{1}=v(p_1^{1},m)$. Em seguida, calcule a \emph{variação compensadora} (CV) do imposto --- o montante que, aos preços com imposto, devolveria o consumidor ao bem-estar pré-imposto --- e interprete o número.
\resolucaobox{7.2cm}
\begin{solucao}
\textbf{Passo 1 --- utilidade indireta.} Substituindo as demandas em $u$:
\[
v(p_1,m)=\left(\dfrac{m}{2p_1}\right)^{1/2}\left(\dfrac{m}{2}\right)^{1/2}=\dfrac{m}{2\sqrt{p_1}}
\] (com $p_2=1$). Em forma geral, $v=\dfrac{m}{2\sqrt{p_1p_2}}$.

\textbf{Passo 2 --- função dispêndio (uma inversão).} Fixando $v=\bar u$ e isolando $m=e$:
\[
\bar u=\frac{e}{2\sqrt{p_1}}\;\Rightarrow\; e(p_1,\bar u)=2\bar u\sqrt{p_1}.
\]
Confirmação por Shephard: $\dfrac{\partial e}{\partial p_1}=2\bar u\cdot\dfrac{1}{2\sqrt{p_1}}=\dfrac{\bar u}{\sqrt{p_1}}=h_1(p_1,\bar u)$ --- consistente.

\textbf{Passo 3 --- níveis de utilidade.}
\[
\bar u^{0}=v(1,72)=\frac{72}{2\sqrt{1}}=36,\qquad \bar u^{1}=v(9,72)=\frac{72}{2\sqrt{9}}=\frac{72}{6}=12.
\]
O imposto corta o bem-estar de $36$ para $12$.

\textbf{Passo 4 --- variação compensadora.} A CV usa a utilidade \emph{inicial} $\bar u^{0}=36$: é a renda extra que, aos preços \emph{com} imposto, restaura esse bem-estar.
\[
\mathrm{CV}=e(p_1^{1},\bar u^{0})-e(p_1^{0},\bar u^{0})=2\cdot 36\cdot\sqrt{9}-2\cdot 36\cdot\sqrt{1}=216-72=144.
\]
Equivalentemente, $e(p_1^{0},\bar u^{0})=m=72$ e $e(p_1^{1},\bar u^{0})=216$, logo $\mathrm{CV}=144$.

\textbf{Intuição econômica:} o imposto custa ao consumidor, em dinheiro de hoje (preços com imposto), R\$~144 de bem-estar --- seria preciso devolver-lhe R\$~144 para recolocá-lo na curva de indiferença original. Esse número é a face dual do problema: em vez de medir a perda em ``utilidade'' (de $36$ para $12$), a função dispêndio a traduz em reais. Repare que a CV $(144)$ supera de longe a receita arrecadada $(32)$ --- a diferença entre o que o consumidor perde e o que o governo ganha é a semente do peso-morto, explorada no próximo item.
\rubricalabel
\begin{itemize}
\item Obtém $v(p_1,m)=m/(2\sqrt{p_1})$: 1 pt.
\item Inverte e obtém $e(p_1,\bar u)=2\bar u\sqrt{p_1}$ (aceita confirmação via Shephard $\partial e/\partial p_1=h_1$): 1 pt.
\item Calcula $\bar u^{0}=36$ e $\bar u^{1}=12$: 1 pt (0{,}5 cada).
\item Calcula $\mathrm{CV}=144$ usando a utilidade inicial $\bar u^{0}=36$: 1{,}5 pt.
\item Interpreta a CV como perda monetária de bem-estar (e observa que excede a receita): 0{,}5 pt.
\item Usar a utilidade \emph{final} $\bar u^{1}$ na CV (confundindo CV com EV): perde os 1{,}5 pt da CV, mantém o restante.
\item Erro aritmético isolado com método correto: penalidade leve de $-0{,}5$ pt.
\end{itemize}
\end{solucao}

\textbf{c)} (5 pontos) O governo cogita substituir o imposto específico por um imposto \emph{lump-sum} (montante fixo) cobrado aos preços não distorcidos ($p_1^{0}=1$). (i) Calcule o lump-sum $L^{*}$ que deixaria o consumidor \emph{exatamente} no mesmo bem-estar a que o imposto específico o levou ($\bar u^{1}=12$); mostre que $L^{*}$ coincide com a \emph{variação equivalente} (EV) do imposto específico. (ii) Compare $L^{*}$ com a receita $R_{\text{arr}}=32$ arrecadada pelo imposto específico e identifique o \emph{peso-morto} (\emph{excess burden}). (iii) Argumente, via dualidade, \emph{por que} o imposto específico é dominado pelo lump-sum equivalente --- isto é, de onde surge a perda de eficiência.
\resolucaobox{7.2cm}
\begin{solucao}
\textbf{Passo 1 --- o lump-sum que iguala o bem-estar.} Um lump-sum $L$ cobrado a preços não distorcidos deixa o consumidor com renda $m-L$ e preço $p_1^{0}=1$, gerando utilidade $v(1,m-L)=\dfrac{m-L}{2}$. Igualando a $\bar u^{1}=12$:
\[
\frac{72-L^{*}}{2}=12\;\Rightarrow\; 72-L^{*}=24\;\Rightarrow\; L^{*}=48.
\]
Um lump-sum de R\$~48 leva o consumidor ao mesmo bem-estar ($\bar u=12$) que o imposto específico.

\textbf{Passo 2 --- $L^{*}$ é a variação equivalente.} A EV mede a perda do imposto aos preços \emph{iniciais}, usando a utilidade \emph{final} $\bar u^{1}=12$:
\[
\mathrm{EV}=e(p_1^{0},\bar u^{0})-e(p_1^{0},\bar u^{1})=2\cdot 36\cdot\sqrt{1}-2\cdot 12\cdot\sqrt{1}=72-24=48=L^{*}.
\]
A coincidência não é acidente: o lump-sum equivalente em bem-estar é, por construção, a quantia que, retirada a preços não distorcidos, reproduz a perda de utilidade do imposto específico --- exatamente a definição de EV.

\textbf{Passo 3 --- peso-morto.} O imposto específico arrecadou $R_{\text{arr}}=32$, mas inflige ao consumidor uma perda equivalente a um lump-sum de $L^{*}=48$. O \emph{excess burden} é o excedente da perda sobre a arrecadação:
\[
\text{peso-morto}=L^{*}-R_{\text{arr}}=\mathrm{EV}-R_{\text{arr}}=48-32=16.
\]
O governo poderia ter extraído R\$~48 via lump-sum causando \emph{a mesma} insatisfação; ao usar o imposto específico, machucou o consumidor o equivalente a $48$ e só arrecadou $32$. Os R\$~16 de diferença evaporam --- não vão para ninguém.

\textbf{Passo 4 --- por que o imposto específico é dominado (dualidade).} A cesta escolhida sob o imposto específico, $x^{t}=(4,36)$, custa, aos preços \emph{verdadeiros} $(1,1)$, exatamente $4+36=40=m-R_{\text{arr}}=72-32$: é portanto \emph{viável} sob o lump-sum de $32$. Sob esse lump-sum, porém, o consumidor \emph{reotimiza} a preços não distorcidos e atinge $v(1,40)=20>12$ --- estritamente melhor que sob o imposto específico. O imposto específico o aprisionou numa cesta com pouco chocolate ($4$ barras) não porque ele ficou pobre, mas porque o preço relativo foi \emph{artificialmente} distorcido: o efeito-substituição do imposto empurrou a escolha para longe do que os custos reais recomendariam. É essa distorção de preço relativo --- ausente no lump-sum --- que gera o peso-morto.

\textbf{Intuição econômica:} o peso-morto é o preço da mentira sobre preços. O lump-sum tira renda e deixa o consumidor alocar livremente a preços que refletem custos reais; o imposto específico tira renda \emph{e} mente sobre o preço relativo do chocolate, induzindo substituição ineficiente. Por isso, para um mesmo sacrifício de bem-estar ($\bar u=12$), o lump-sum arrecada mais ($48$ contra $32$). A lição vale tanto para o desenho de qualquer imposto sobre o consumo quanto, mais geralmente, para a escolha entre tributos distorcivos e não distorcivos --- tema que retorna em tributação e em equilíbrio geral com impostos.
\rubricalabel
\begin{itemize}
\item Calcula $L^{*}=48$ via $v(1,m-L^{*})=\bar u^{1}=12$: 1{,}5 pt.
\item Mostra que $L^{*}=\mathrm{EV}=e(p_1^{0},\bar u^{0})-e(p_1^{0},\bar u^{1})=48$: 1 pt.
\item Identifica o peso-morto $=L^{*}-R_{\text{arr}}=48-32=16$: 1 pt.
\item Argumenta via dualidade por que o imposto específico é dominado --- viabilidade da cesta $x^{t}$ sob o lump-sum e reotimização a preços não distorcidos / distorção do preço relativo como fonte do peso-morto: 1{,}5 pt.
\item Apenas afirmar que o lump-sum é melhor sem justificativa estrutural (substituição/dualidade): teto de 3 pts no item.
\item Erro aritmético isolado com método correto (ex.: monta $L^{*}$ certo, erra a subtração final): penalidade leve de $-0{,}5$ pt.
\end{itemize}
\end{solucao}



\section*{Bloco 2 --- Equilíbrio Geral \hfill \normalsize (25 pontos)}

\textbf{Questão 3.} (10 pontos) \textbf{Itens objetivos.} Nos itens \textbf{a)} e \textbf{b)}, assinale a \emph{única} alternativa correta (3 pontos cada; não há penalização por erro). Nos itens \textbf{c)} e \textbf{d)}, responda \textbf{Verdadeiro ou Falso} (2 pontos cada; vale a regra de anulação descrita na capa). Registre suas respostas no quadro-resposta ao final do bloco.

\textbf{a)} (3 pontos) A conclusão das negociações Mercosul--União Europeia (dezembro de 2024) reacendeu a discussão sobre a troca de grãos (bem $1$) por manufaturas (bem $2$). Modele o comércio birregional como uma economia de trocas pura $2\times 2$, com o Brasil ($A$) e a União Europeia ($B$) como agentes. Suponha que as preferências sejam $u^A(x_1,x_2)=x_1^{3/4}x_2^{1/4}$ e $u^B(x_1,x_2)=x_1^{1/2}x_2^{1/2}$, e que as dotações iniciais sejam $\omega^A=(8,0)$ (o Brasil chega com grãos) e $\omega^B=(0,8)$ (a UE chega com manufaturas). Tomando o bem $2$ como numerário ($p_2=1$), o preço relativo de equilíbrio walrasiano dos grãos, $p_1^{*}$, é igual a:

\begin{enumerate}[label=\Alph*., leftmargin=3.2em, itemsep=0.15em, topsep=0.2em]
  \item $p_1^{*}=2$
  \item $p_1^{*}=\dfrac{1}{2}$
  \item $p_1^{*}=3$
  \item $p_1^{*}=\dfrac{3}{2}$
  \item $p_1^{*}=1$
\end{enumerate}
\gabaritoitem{A}
\begin{solucao}
{\footnotesize\textit{Paralelo de: Simulado 1, Bloco 2, ME1 (equilíbrio walrasiano de trocas Cobb-Douglas, Brasil–China soja$\times$manufaturas).}}\par\vspace{0.2em}
\textbf{Passo 1 --- O que é dado e o que se pede.} Dois agentes Cobb-Douglas, dotações nos cantos, $p_2=1$. Pede-se o preço relativo $p_1^{*}$ que zera o excesso de demanda agregada pelos grãos.

\textbf{Passo 2 --- Plano.} Para Cobb-Douglas $x_1^{\alpha}x_2^{1-\alpha}$, a demanda marshalliana pelo bem $1$ é $x_1(p,m)=\alpha\,m/p_1$. Basta escrever as rendas $m^i=p\cdot\omega^i$, somar as demandas pelo bem $1$ e igualar à oferta total $\bar\omega_1=8$ (pela Lei de Walras, equilibrar um mercado equilibra o outro).

\textbf{Passo 3 --- Execução.} As rendas são $m^A=p_1\cdot 8+1\cdot 0=8p_1$ e $m^B=p_1\cdot 0+1\cdot 8=8$. Com $\alpha^A=\tfrac34$ e $\alpha^B=\tfrac12$:
\[ x_1^A=\tfrac34\cdot\frac{8p_1}{p_1}=6, \qquad x_1^B=\tfrac12\cdot\frac{8}{p_1}=\frac{4}{p_1}. \]
Equilíbrio no mercado de grãos: $x_1^A+x_1^B=\bar\omega_1$, ou seja, $6+\dfrac{4}{p_1}=8\Rightarrow\dfrac{4}{p_1}=2\Rightarrow p_1^{*}=2$.

\textbf{Passo 4 --- Verificação.} No mercado de manufaturas: $x_2^A=\tfrac14\cdot 8p_1=2p_1=4$ e $x_2^B=\tfrac12\cdot 8=4$; soma $=8=\bar\omega_2$. Os dois mercados fecham, confirmando a Lei de Walras.

\textbf{Por que os distratores mais tentadores enganam.} A alternativa (B), $p_1^{*}=\tfrac12$, é a inversão da razão de preços: quem lê $p_1$ como ``manufaturas por grãos'' (em vez de ``grãos em unidades de manufatura'') reporta o recíproco. A alternativa (E), $p_1^{*}=1$, surge ao impor simetria indevida $\alpha^A=\alpha^B=\tfrac12$ (esquecendo que o Brasil pondera $\tfrac34$ os grãos): aí $x_1^A=4$, $4+4/p_1=8\Rightarrow p_1=1$. Como o Brasil traz os grãos (bem abundante) mas também os valoriza intensamente, a demanda interna sustenta um preço relativo alto: cada unidade de grão vale duas de manufatura no equilíbrio.

\textbf{Intuição econômica:} o preço relativo é o mensageiro que reconcilia ofertas e desejos. Aqui, embora o Brasil chegue ao mercado com os grãos, sua própria preferência forte por esse bem ($\alpha^A=\tfrac34$) mantém a demanda agregada elevada, e o equilíbrio precifica os grãos em duas unidades de manufatura --- abundância de dotação não basta para baratear um bem se quem o detém também o deseja muito.
\end{solucao}

\textbf{b)} (3 pontos) O complexo termelétrico da GNA (Gás Natural Açu), no Porto do Açu (RJ), é o maior parque a gás natural da América Latina, com receita que depende do despacho do operador do sistema. Há dois estados da natureza para o próximo ciclo: $s_1$ = térmica despachada (receita alta) e $s_2$ = fora da ordem de despacho (receita baixa). O mercado é completo e livre de arbitragem. Negociam-se: o título $X$ --- paga R\$~4 em $s_1$, R\$~2 em $s_2$, custa R\$~3{,}00 --- e o título $Y$ --- paga R\$~1 em $s_1$, R\$~3 em $s_2$, custa R\$~1{,}50. A GNA quer precificar um contrato $Z$ que paga R\$~20 em $s_1$ e R\$~30 em $s_2$. Pela ausência de arbitragem, o preço de $Z$ é:

\begin{enumerate}[label=\Alph*., leftmargin=3.2em, itemsep=0.15em, topsep=0.2em]
  \item R\$~15
  \item R\$~24
  \item R\$~21
  \item R\$~30
  \item R\$~9
\end{enumerate}
\gabaritoitem{C}
\begin{solucao}
{\footnotesize\textit{Paralelo de: Simulado 3, Bloco 2, ME2 (precificação Arrow-Debreu por não-arbitragem, Serena Energia/eólica).}}\par\vspace{0.2em}
\textbf{Passo 1 --- O que é dado e o que se pede.} Mercado completo com dois estados; dois títulos $X$ e $Y$ com \emph{payoffs} e preços conhecidos. Pede-se o preço de um terceiro contrato $Z$ por não-arbitragem --- isto é, pelo custo de \emph{replicá-lo} com $X$ e $Y$, ou, equivalentemente, pelos preços de estado (Arrow) implícitos.

\textbf{Passo 2 --- Recuperar os preços de Arrow.} Sejam $q_1,q_2$ os preços dos títulos elementares que pagam R\$~1 num estado e $0$ no outro. Os preços observados impõem o sistema
\[ 4q_1+2q_2=3{,}00, \qquad q_1+3q_2=1{,}50. \]
Multiplicando a segunda por $4$: $4q_1+12q_2=6{,}00$; subtraindo a primeira: $10q_2=3{,}00\Rightarrow q_2=0{,}30$. Então $q_1=1{,}50-3\cdot 0{,}30=1{,}50-0{,}90=0{,}60$. Os preços de estado são únicos porque o mercado é completo.

\textbf{Passo 3 --- Precificar $Z$.} Por não-arbitragem, o preço de qualquer \emph{payoff} $(z_1,z_2)$ é $q_1 z_1+q_2 z_2$:
\[ P_Z=0{,}60\cdot 20+0{,}30\cdot 30=12+9=21. \]
Logo $P_Z=$ R\$~21 --- alternativa (C).

\textbf{Passo 4 --- Verificação.} A carteira que replica $Z$ deve pagar $(20,30)$: resolvendo $4a+b=20$, $2a+3b=30$ em unidades de $X$ e $Y$ chega-se a $a=3$, $b=8$, cujo custo é $3\cdot 3{,}00+8\cdot 1{,}50=9{,}00+12{,}00=21$. O custo de replicação coincide com $P_Z$ --- confirmação da não-arbitragem.

\textbf{Por que os distratores mais tentadores enganam.} A alternativa (B), R\$~24, troca os \emph{payoffs} de estado: precifica $(30,20)$ em vez de $(20,30)$, $0{,}60\cdot 30+0{,}30\cdot 20$ --- erro de associar o pagamento maior ao estado errado. A alternativa (D), R\$~30, toma ambos os preços de estado iguais a $q_1=0{,}60$ ($0{,}60\cdot 50$), ignorando que os preços de Arrow são \emph{desiguais}. A alternativa (A), R\$~15, faz o oposto ($q_1=q_2=0{,}30$); a (E), R\$~9, conta só a perna de $s_2$ ($0{,}30\cdot 30$).

\textbf{Intuição econômica:} num mercado completo e sem arbitragem, todo contrato contingente vale o que custa montá-lo com os tijolos elementares --- os títulos de Arrow de cada estado. A GNA não precisa adivinhar probabilidades de despacho: os preços de mercado dos contratos negociados já embutem o valor de R\$~1 entregue com a térmica despachada ($0{,}60$) e fora de ordem ($0{,}30$). O contrato $Z$ vale a soma ponderada desses preços de estado: R\$~21.
\end{solucao}

\textbf{c)} (2 pontos, V ou F) Com a privatização da Eletrobras concluída em 2022, leilões de capacidade de geração passaram a alocar contratos entre geradoras que tomam os preços como dados. Modele um desses ambientes como uma economia competitiva de trocas com preferências contínuas e localmente não-saciadas, mas \emph{não} necessariamente convexas. \textbf{Afirmação:} o $2^{\text{o}}$ Teorema do Bem-Estar --- toda alocação Pareto-eficiente pode ser sustentada como equilíbrio walrasiano com transferências \emph{lump-sum} apropriadas --- vale sob exatamente as mesmas hipóteses do $1^{\text{o}}$ Teorema (continuidade e não-saciedade local dos agentes), \emph{sem} que seja necessário supor convexidade das preferências.

\gabaritoitem{F}
\begin{solucao}
{\footnotesize\textit{Paralelo de: Simulado 2, Bloco 2, VF1 (papel da convexidade nos teoremas do bem-estar; gabarito F; cooperativas de café de Minas Gerais).}}\par\vspace{0.2em}
\textbf{Passo 1 --- O que se afirma.} Que o $2^{\text{o}}$ Teorema do Bem-Estar (Pareto $\Rightarrow$ descentralizável por preços mais transferências) vale com as \emph{mesmas} hipóteses do $1^{\text{o}}$ --- ou seja, \emph{sem} convexidade das preferências.

\textbf{Passo 2 --- O que cada teorema realmente exige.} O $1^{\text{o}}$ Teorema (equilíbrio competitivo $\Rightarrow$ Pareto-eficiência) é puramente orçamentário e requer apenas não-saciedade local: se uma alocação viável Pareto-dominasse o equilíbrio, cada cesta fracamente preferida custaria pelo menos a renda do agente (e a estritamente preferida, estritamente mais); somando as restrições, o valor agregado da alocação dominante excederia o da dotação --- contradição com a viabilidade. Convexidade nunca é usada. O $2^{\text{o}}$ Teorema, ao contrário, parte de uma alocação eficiente e \emph{constrói} um vetor de preços que a sustenta; essa construção apoia-se no Teorema da Separação de Hiperplanos aplicado à soma dos conjuntos ``estritamente preferidos''. Esse conjunto-soma só é garantidamente convexo se as preferências individuais forem convexas.

\textbf{Por que é Falsa.} Sem convexidade, o conjunto agregado de cestas preferidas pode ter ``reentrâncias'', e não existe hiperplano de preços que o separe da dotação eficiente: a alocação é Pareto-eficiente mas \emph{não} descentralizável por nenhum sistema de preços lineares. A convexidade é, portanto, hipótese \emph{adicional indispensável} do $2^{\text{o}}$ Teorema --- ele \emph{não} vale sob as mesmas hipóteses do $1^{\text{o}}$, e a afirmação é falsa.

\textbf{Confusão entre os dois teoremas.} A afirmação equipara o teorema fácil ($1^{\text{o}}$, robusto à forma das preferências) ao teorema difícil ($2^{\text{o}}$, que paga o pedágio da convexidade). São níveis distintos de exigência: o $1^{\text{o}}$ certifica que um equilíbrio que ocorre é eficiente; o $2^{\text{o}}$ promete realizar qualquer eficiência desejada via preços, e essa promessa precisa de geometria convexa.

\textbf{Intuição econômica:} a mão invisível ``para frente'' (mercado $\Rightarrow$ eficiência) funciona contabilmente, via orçamentos, e quase não cobra hipóteses sobre o formato das preferências. A mão invisível ``para trás'' (escolher uma alocação eficiente e implementá-la com preços e transferências) precisa separar conjuntos preferidos por um hiperplano de preços --- e separar exige convexidade. No leilão de capacidade, se uma geradora tem preferências não-convexas, um \emph{planner} pode até identificar a alocação eficiente desejada, mas talvez não consiga sustentá-la com preço algum de \emph{clearing} mais transferências.
\end{solucao}

\textbf{d)} (2 pontos, V ou F) A WEG, fabricante catarinense de motores elétricos e grande exportadora, tem fluxo de caixa futuro sensível à taxa de câmbio. Modele o próximo trimestre por dois estados da natureza: $s_1$ = real apreciado (receita de exportação em reais menor) e $s_2$ = real depreciado (receita maior). Suponha que o \emph{único} ativo negociável seja um título que paga $2$ em $s_1$ e $1$ em $s_2$, e que nenhum outro contrato contingente exista. \textbf{Afirmação:} mesmo sendo um ativo cujo retorno \emph{varia} entre os estados, ele não basta para transferir recursos do estado favorável ($s_2$) para o desfavorável ($s_1$): o espaço de carteiras atingíveis é unidimensional, de modo que o risco cambial específico permanece não-segurável.

\gabaritoitem{V}
\begin{solucao}
{\footnotesize\textit{Paralelo de: Simulado 3, Bloco 2, VF1 (mercado incompleto com único ativo negociável; gabarito V; Embraer/tarifaço).}}\par\vspace{0.2em}
\textbf{Passo 1 --- O que se afirma.} Que, mesmo com um ativo cujo \emph{payoff} difere entre estados --- $(2,1)$ --- mas sendo ele o \emph{único} negociável, o risco \emph{estado-específico} (caixa baixo sob $s_1$, alto sob $s_2$) ainda não pode ser transferido: o mercado é incompleto e o risco, não-segurável.

\textbf{Passo 2 --- O espaço gerado pelo ativo.} Comprar $\theta$ unidades do título gera o \emph{payoff} $\theta\cdot(2,1)=(2\theta,\theta)$. O conjunto de \emph{payoffs} atingíveis é a reta $\{(2\theta,\theta):\theta\in\mathbb{R}\}$ --- um subespaço de dimensão $1$ dentro de $\mathbb{R}^2$. O ponto crucial é que a incompletude \emph{não} vem de o ativo ``pagar o mesmo nos dois estados'' (ele não paga: $2\neq 1$), e sim de existir um \emph{único} ativo linearmente independente para $2$ estados.

\textbf{Passo 3 --- Por que o risco não é transferível.} Transferir recursos do estado favorável $s_2$ para o desfavorável $s_1$ exigiria um \emph{payoff} líquido com \emph{sinais opostos} entre estados, da forma $(+,-)$. Mas todo elemento do \emph{span}, $(2\theta,\theta)$, tem as duas componentes \emph{com o mesmo sinal} de $\theta$ (e na razão fixa $2:1$); nenhum \emph{payoff} $(a,b)$ com $a$ e $b$ de sinais contrários é replicável. Como há $2$ estados e apenas $1$ ativo independente, e $1<2$, o mercado é \emph{incompleto}. A afirmação é \textbf{verdadeira}.

\textbf{Passo 4 --- Armadilha.} O aluno desatento pensa: ``o ativo varia entre estados, logo permite \emph{hedge}''. Mas variar entre estados \emph{não} é o mesmo que permitir transferir risco \emph{em qualquer direção}. O título $(2,1)$ desloca consumo na proporção rígida $2:1$ (sempre mais em $s_1$ quando $\theta>0$); ele jamais entrega o perfil $(+,-)$ necessário para compensar especificamente o estado ruim sem mexer no bom. Para isso seria preciso um \emph{segundo} ativo linearmente independente.

\textbf{Intuição econômica:} completar mercados é ter tantos ativos \emph{linearmente independentes} quantos forem os estados da natureza --- não basta que o único ativo disponível seja ``arriscado''. Um só papel, por mais que oscile entre real apreciado e depreciado, traça apenas uma reta de \emph{payoffs}: cobre a direção dele, deixa todo o resto do plano de riscos a descoberto. Para se proteger do câmbio, a WEG precisaria de um instrumento adicional --- por exemplo, um contrato cambial a termo --- que, combinado ao título existente, gerasse perfis $(+,-)$; sem ele, o risco específico permanece no balanço.
\end{solucao}

\quadroresposta{2}

\textbf{Questão 4.} (15 pontos) A safra brasileira de soja 2024/2025 bateu recorde, mas o ciclo seguinte traz incerteza climática associada ao fenômeno La Niña, que ameaça a produtividade no Sul. Considere uma cooperativa de produtores de soja cujo fluxo de caixa do próximo ano depende do regime de chuvas. Modele o futuro como dois estados da natureza mutuamente exclusivos: $s_1$ = clima favorável (safra cheia, caixa alto) e $s_2$ = La Niña (seca, quebra de safra, caixa baixo). Suponha probabilidades $\pi_1=\tfrac23$ e $\pi_2=\tfrac13$. A cooperativa consome um único bem físico em cada estado e tem dotação contingente de fluxo de caixa $\omega=(\omega_1,\omega_2)=(90,30)$ (ganha mais quando o clima coopera). Suas preferências são de utilidade esperada com Bernoulli $v(c)=\ln c$ e sem desconto temporal. Existe um mercado completo de bens contingentes (Arrow-Debreu) onde se negociam títulos que pagam $1$ unidade do bem em um estado e $0$ no outro, a preços $q=(q_1,q_2)$.

\textbf{a)} (5 pontos) Suponha que os preços dos títulos contingentes sejam \emph{atuarialmente justos}, isto é, $q_s=\pi_s$, de modo que $q=\left(\tfrac23,\tfrac13\right)$. Monte o problema de maximização de utilidade esperada da cooperativa sujeito à restrição orçamentária contingente, derive a condição de primeira ordem e determine o consumo contingente ótimo $c^{*}=(c_1^{*},c_2^{*})$. Interprete o formato da solução.
\resolucaobox{7.2cm}
\begin{solucao}
{\footnotesize\textit{Paralelo de: Simulado 1, Bloco 2, Aberta (bens contingentes Arrow-Debreu com utilidade esperada ln; geradora hidrelétrica/ONS).}}\par\vspace{0.2em}
\textbf{Passo 1 --- Montagem.} A riqueza inicial avaliada a mercado é $W=q\cdot\omega=\tfrac23\cdot 90+\tfrac13\cdot 30=60+10=70$. O problema é
\[ \max_{c_1,c_2}\;\pi_1\ln c_1+\pi_2\ln c_2\quad\text{s.a.}\quad q_1 c_1+q_2 c_2=W. \]

\textbf{Passo 2 --- Lagrangiano e CPO.} Com $\mathcal{L}=\pi_1\ln c_1+\pi_2\ln c_2+\lambda\,(W-q_1 c_1-q_2 c_2)$:
\[ \frac{\partial\mathcal{L}}{\partial c_s}=\frac{\pi_s}{c_s}-\lambda q_s=0\;\Longrightarrow\; c_s=\frac{\pi_s}{\lambda\,q_s}. \]

\textbf{Passo 3 --- Preços justos.} Com $q_s=\pi_s$, cada $c_s=\dfrac{\pi_s}{\lambda\,\pi_s}=\dfrac{1}{\lambda}$, \emph{constante entre estados}. Logo $c_1^{*}=c_2^{*}$. Substituindo na restrição: $q_1 c+q_2 c=(\pi_1+\pi_2)c=c=W=70$. Portanto $c^{*}=(70,70)$.

\textbf{Passo 4 --- Verificação.} A operação de seguro é autofinanciada: a cooperativa vende $90-70=20$ do título de clima favorável e compra $70-30=40$ do título de seca; o caixa líquido é $q_2\cdot 40-q_1\cdot 20=\tfrac13\cdot 40-\tfrac23\cdot 20=\tfrac{40}{3}-\tfrac{40}{3}=0$, coerente com $W$ inalterada. Seguro completo a preço justo não altera a riqueza, só sua distribuição entre estados.

\textbf{Intuição econômica:} sob aversão ao risco ($v$ côncava) e preços atuarialmente justos, o agente compra \emph{seguro completo} e iguala o consumo entre os estados ($c_1=c_2$) --- a clássica regra de seguro total. A cooperativa troca o excedente de caixa dos anos de safra cheia por proteção nos anos de La Niña, suavizando perfeitamente seu fluxo real e tirando o risco climático do balanço.
\rubricalabel
\begin{itemize}
\item Calcula a riqueza $W=q\cdot\omega=70$ corretamente: 1 pt.
\item Monta o problema de maximização de utilidade esperada com a restrição orçamentária contingente correta: 1 pt.
\item Deriva a CPO $\pi_s/c_s=\lambda q_s$ (lagrangiano ou substituição equivalente): 1{,}5 pt.
\item Conclui $c_1^{*}=c_2^{*}$ sob preços justos e obtém $c^{*}=(70,70)$: 1 pt.
\item Interpreta como seguro completo / suavização de consumo: 0{,}5 pt.
\item Erro aritmético com método correto (CPO certa, mas conta de $W$ ou substituição final trocada): penalidade leve de no máximo $0{,}5$ pt.
\end{itemize}
\end{solucao}

\textbf{b)} (5 pontos) Agora suponha que o mercado precifique os títulos contingentes em $q=\left(\tfrac12,\tfrac12\right)$ --- ou seja, o título que paga na \emph{seca} ficou caro relativamente à sua probabilidade. Recalcule o consumo contingente ótimo $c^{*}=(c_1^{*},c_2^{*})$ e compare com o item anterior, explicando por que a cooperativa deixa de se segurar plenamente.
\resolucaobox{7.2cm}
\begin{solucao}
\textbf{Passo 1 --- Nova riqueza.} $W=q\cdot\omega=\tfrac12\cdot 90+\tfrac12\cdot 30=60$.

\textbf{Passo 2 --- Razão de consumo via CPO.} Da CPO $c_s=\pi_s/(\lambda q_s)$, a razão entre estados é
\[ \frac{c_1}{c_2}=\frac{\pi_1/q_1}{\pi_2/q_2}=\frac{\pi_1\,q_2}{\pi_2\,q_1}=\frac{\tfrac23\cdot\tfrac12}{\tfrac13\cdot\tfrac12}=\frac{\tfrac23}{\tfrac13}=2. \]
Logo $c_1=2c_2$ (consome \emph{mais} no estado de clima favorável).

\textbf{Passo 3 --- Resolver a restrição.} $q_1 c_1+q_2 c_2=\tfrac12(2c_2)+\tfrac12 c_2=\tfrac32 c_2=60\Rightarrow c_2^{*}=40$ e $c_1^{*}=80$. Assim $c^{*}=(80,40)$.

\textbf{Passo 4 --- Verificação e armadilha.} Custo total: $\tfrac12\cdot 80+\tfrac12\cdot 40=40+20=60=W$. A armadilha é supor que aversão ao risco \emph{sempre} implica $c_1=c_2$: isso só vale com preços justos. Aqui o seguro da seca está \emph{carregado} (preço $q_2=\tfrac12$ contra probabilidade $\pi_2=\tfrac13$), então a cooperativa compra menos proteção e aceita carregar parte do risco climático.

\textbf{Intuição econômica:} a regra geral é igualar a taxa marginal de substituição contingente à razão de preços: $\dfrac{\pi_1 v'(c_1)}{\pi_2 v'(c_2)}=\dfrac{q_1}{q_2}$. Quando o preço do título de seca excede sua probabilidade, segurar-se contra a La Niña fica caro na margem, e o agente racionalmente \emph{sub-segura} esse estado --- consome mais na safra cheia e menos na seca. A suavização perfeita do item (a) era um caso especial dos preços justos, não uma lei universal da aversão ao risco.
\rubricalabel
\begin{itemize}
\item Recalcula a riqueza $W=60$ com os novos preços: 1 pt.
\item Obtém a razão $c_1/c_2=(\pi_1 q_2)/(\pi_2 q_1)=2$ a partir da CPO: 1{,}5 pt.
\item Resolve a restrição e chega a $c^{*}=(80,40)$: 1{,}5 pt.
\item Explica corretamente por que não há seguro completo (preço da seca acima do atuarial / título carregado): 1 pt.
\item Erro aritmético com método correto (razão certa mas conta final trocada): penalidade leve de no máximo $0{,}5$ pt.
\end{itemize}
\end{solucao}

\textbf{c)} (5 pontos) Mantenha os preços do item (b), $q=\left(\tfrac12,\tfrac12\right)$. A cooperativa deseja contratar um \emph{seguro contra a La Niña}: um título que paga $1$ no estado $s_2$ (seca) e $0$ no estado $s_1$ (clima favorável). (i) Pela ausência de arbitragem em mercado completo, qual é o preço desse seguro? Compare-o com o prêmio atuarialmente justo $\pi_2$ e diga se ele está barato ou \emph{carregado}. (ii) Suponha agora que apenas o título de \emph{clima favorável} (que paga $1$ em $s_1$ e $0$ em $s_2$) seja negociável, e nenhum outro ativo. Classifique este segundo mercado como \emph{completo} ou \emph{incompleto} e justifique em uma frase, indicando se o seguro contra a seca é replicável.
\resolucaobox{7.2cm}
\begin{solucao}
\textbf{Passo 1 --- Preço do seguro contra a seca (mercado completo).} O seguro contra a La Niña é exatamente o título de Arrow do estado $s_2$, com \emph{payoff} $(0,1)$. Sob ausência de arbitragem, seu preço é o próprio preço de Arrow desse estado:
\[ P=q_2=\tfrac12. \]

\textbf{Passo 2 --- Comparação com o prêmio justo.} O prêmio atuarialmente justo seria $\pi_2=\tfrac13$. Como $P=\tfrac12>\tfrac13=\pi_2$, o seguro está \emph{carregado}: o \emph{loading} é $\dfrac{q_2}{\pi_2}-1=\dfrac{1/2}{1/3}-1=\dfrac32-1=\tfrac12$, ou seja, $50\%$ acima do justo. Isso é coerente com o item (b), em que a cooperativa, vendo o seguro caro, escolheu \emph{sub-segurar} a seca.

\textbf{Passo 3 --- Classificação do segundo mercado.} Com apenas o título de clima favorável negociável, o único \emph{payoff} replicável é da forma $x\cdot(1,0)=(x,0)$ --- o espaço gerado (\emph{span}) é a reta $\{(x,0):x\in\mathbb{R}\}$, de dimensão $1$. Como $1<2=|S|$ (número de estados), o mercado é \emph{incompleto}. Em particular, o seguro contra a seca $(0,1)$ \emph{não} pertence a esse \emph{span}: \emph{não é replicável}, e a cooperativa fica com risco de La Niña não-segurável.

\textbf{Passo 4 --- Verificação.} Coerência interna: no mercado completo o seguro contra a seca \emph{é} o título de Arrow de $s_2$, então seu preço tem de ser $q_2$ --- e de fato $P=q_2=\tfrac12$. Já no mercado com um só título, falta um estado independente ($1$ ativo para $2$ estados), o que torna impossível replicar qualquer \emph{payoff} com componente em $s_2$.

\textbf{Intuição econômica:} um mercado é completo quando há tantos ativos linearmente independentes quanto estados da natureza --- só então todo risco contingente pode ser comprado ou vendido. Retirar um dos títulos de Arrow abre um ``buraco'' no espaço de riscos: a cooperativa deixa de conseguir transferir recursos para o estado de seca, e nenhum preço de mercado consegue eliminar essa exposição. A completude é, portanto, a condição que dá sentido a falar em \emph{preço único} de qualquer seguro.
\rubricalabel
\begin{itemize}
\item Identifica o seguro contra a seca como o título de Arrow de $s_2$ e o precifica em $P=q_2=\tfrac12$ por não-arbitragem: 1{,}5 pt.
\item Compara com o justo $\pi_2=\tfrac13$ e conclui que está carregado (\emph{loading} de $50\%$): 1 pt.
\item Classifica corretamente o mercado com um só título como incompleto, com justificativa de \emph{span}/dimensão $1<2=|S|$: 1{,}5 pt.
\item Conclui que o seguro contra a seca não é replicável nesse mercado incompleto: 1 pt.
\item Erro aritmético com método correto (não-arbitragem certa, conta de \emph{loading} trocada): penalidade leve de no máximo $0{,}5$ pt.
\end{itemize}
\end{solucao}



\section*{Bloco 3 --- Externalidades e Bens Públicos \hfill \normalsize (25 pontos)}

\textbf{Questão 5.} (10 pontos) \textbf{Itens objetivos.} Nos itens \textbf{a)} e \textbf{b)}, assinale a \emph{única} alternativa correta (3 pontos cada; não há penalização por erro). Nos itens \textbf{c)} e \textbf{d)}, responda \textbf{Verdadeiro ou Falso} (2 pontos cada; vale a regra de anulação descrita na capa). Registre suas respostas no quadro-resposta ao final do bloco.

\textbf{a)} (3 pontos) O M\'etodo Wolbachia (Fiocruz / Minist\'erio da Sa\'ude) libera mosquitos \emph{Aedes aegypti} que n\~ao transmitem a dengue; soltar lotes num bairro tamb\'em reduz a circula\c{c}\~ao do v\'irus nos bairros vizinhos --- uma externalidade positiva. Numa certa cidade, o benef\'icio marginal \emph{privado} de soltar lotes \'e a demanda inversa $P(q)=40-2q$ (em R\$~mil por lote), o custo marginal privado \'e constante, $\text{MPC}=10$, e cada lote gera um benef\'icio externo marginal constante de R\$~8 mil aos vizinhos. Qual \'e o subs\'idio pigouviano \'otimo por lote que conduz a cidade \`a libera\c{c}\~ao socialmente eficiente?

\begin{enumerate}[label=\Alph*., leftmargin=3.2em, itemsep=0.15em, topsep=0.2em]
  \item R\$~2
  \item R\$~4
  \item R\$~10
  \item R\$~8
  \item R\$~16
\end{enumerate}
\gabaritoitem{D}
\begin{solucao}
{\footnotesize\textit{Paralelo de: Simulado 2, Bloco 3, ME1 (subsídio pigouviano / polinização abelhas-laranjais de Limeira).}}\par\vspace{0.2em}
\textbf{Passo 1 --- o que se pede.} Numa externalidade \emph{positiva}, o benef\'icio social de cada lote excede o privado pelo benef\'icio externo marginal ($\text{MEB}$) que recai sobre os bairros vizinhos. Deixada ao c\'alculo privado, a cidade libera \emph{pouco}; o subs\'idio pigouviano \'otimo fecha exatamente essa diferen\c{c}a, igualando $\text{MEB}$ na margem eficiente.

\textbf{Passo 2 --- equil\'ibrio privado.} Deixado sozinho, o agente libera lotes at\'e que seu benef\'icio marginal privado iguale o custo marginal privado:
\[
P(q)=\text{MPC}\;\Longrightarrow\;40-2q=10\;\Longrightarrow\; q^{\text{priv}}=15.
\]

\textbf{Passo 3 --- \'otimo social e subs\'idio.} O \'otimo social iguala o \emph{benef\'icio marginal social} ($P(q)+\text{MEB}$) ao custo marginal:
\[
P(q)+\text{MEB}=\text{MPC}\;\Longrightarrow\;(40-2q)+8=10\;\Longrightarrow\; q^{*}=19.
\]
O subs\'idio pigouviano \'otimo \'e o benef\'icio externo marginal avaliado no \'otimo. Como $\text{MEB}=8$ \'e constante, $s^{*}=8$. Verifica\c{c}\~ao: com o subs\'idio, o agente maximiza $P(q)+s^{*}=\text{MPC}$, isto \'e $40-2q=10-8=2\Rightarrow q=19=q^{*}$. O subs\'idio empurra a libera\c{c}\~ao de $15$ para $19$ lotes.

\textbf{Passo 4 --- distratores.} A alternativa $R\$~10$ confunde o \emph{subs\'idio} com o $\text{MPC}=10$, valor saliente no enunciado --- toma o custo no lugar do diferencial de benef\'icio marginal; \'e o erro mais comum. A alternativa $R\$~4$ corresponde \`a \emph{varia\c{c}\~ao} de quantidade $q^{*}-q^{\text{priv}}=19-15=4$: identifica o subs\'idio com o deslocamento de output, n\~ao com o tamanho da externalidade. A alternativa $R\$~2$ \'e o pre\c{c}o de demanda $P(q^{*})=40-2\cdot 19=2$ no \'otimo (o custo l\'iquido por lote ap\'os o subs\'idio), confundido com o pr\'oprio subs\'idio. A alternativa $R\$~16=2\,\text{MEB}$ \'e erro de escala (dobrar a externalidade).

\textbf{Intui\c{c}\~ao econ\^omica:} cada lote liberado gera um ganho de R\$~8 mil que o agente n\~ao captura --- protege os vizinhos de gra\c{c}a. Por isso libera de menos, ignorando R\$~8 mil de benef\'icio social por lote. O subs\'idio de R\$~8 mil internaliza esse ganho externo, alinhando o incentivo privado ao social. O valor do subs\'idio \'e o tamanho da externalidade na margem --- n\~ao a quantidade-alvo ($q^{*}=19$) nem sua varia\c{c}\~ao ($R\$~4$); o valor $R\$~10$ \'e apenas o $\text{MPC}$, o custo de produzir e liberar cada lote, e n\~ao o tamanho do benef\'icio externo que se quer internalizar.
\end{solucao}

\textbf{b)} (3 pontos) Desde janeiro de 2025, Manhattan opera o primeiro ped\'agio de congestionamento dos Estados Unidos: ve\'iculos que entram na zona central, ao sul da Rua 60, pagam tarifa no hor\'ario de pico. Antes da tarifa, o acesso funcionava como \emph{recurso comum} --- rival (mais carros congestionam todos) mas sem exclus\~ao. Seja $N$ o n\'umero de viagens de autom\'ovel por hora-pico (em milhares), o \emph{benef\'icio bruto total} da via $B(N)=N\,(90-N)$ (em R\$~mil) e o custo privado de cada viagem constante e igual a $30$. Sob acesso livre, o n\'umero de viagens $N^{\text{AL}}$ e o socialmente \'otimo $N^{*}$ s\~ao, respectivamente:

\begin{enumerate}[label=\Alph*., leftmargin=3.2em, itemsep=0.15em, topsep=0.2em]
  \item $N^{\text{AL}}=30,\quad N^{*}=60$
  \item $N^{\text{AL}}=60,\quad N^{*}=30$
  \item $N^{\text{AL}}=60,\quad N^{*}=60$
  \item $N^{\text{AL}}=45,\quad N^{*}=30$
  \item $N^{\text{AL}}=90,\quad N^{*}=45$
\end{enumerate}
\gabaritoitem{B}
\begin{solucao}
{\footnotesize\textit{Paralelo de: Simulado 3, Bloco 3, ME2 (tragédia dos comuns / órbita baixa Starlink — médio vs.}}\par\vspace{0.2em}
\textbf{Passo 1 --- o que se pede.} A trag\'edia dos comuns nasce de uma cunha entre o que o motorista \emph{privado} percebe e o que o planejador deveria considerar. Sob acesso livre, cada motorista embolsa o benef\'icio \emph{m\'edio} da via, $B(N)/N$, e entra enquanto esse m\'edio cobrir o custo; o planejador, ao contr\'ario, compara o benef\'icio \emph{marginal} $B'(N)$ ao custo. A cunha entre m\'edio e marginal \'e a externalidade de congestionamento que o motorista individual n\~ao paga.

\textbf{Passo 2 --- benef\'icio m\'edio e marginal.} De $B(N)=N(90-N)=90N-N^{2}$:
\[
\text{benef\'icio m\'edio}=\frac{B(N)}{N}=90-N,\qquad \text{benef\'icio marginal}=B'(N)=90-2N.
\]

\textbf{Passo 3 --- acesso livre.} Motoristas entram enquanto o benef\'icio m\'edio cobrir o custo; a entrada cessa quando se igualam:
\[
90-N=30\;\Longrightarrow\; N^{\text{AL}}=60.
\]

\textbf{Passo 4 --- \'otimo social.} O planejador iguala o benef\'icio marginal ao custo:
\[
90-2N=30\;\Longrightarrow\;2N=60\;\Longrightarrow\; N^{*}=30.
\]
Logo $N^{\text{AL}}=60>N^{*}=30$: o acesso livre coloca o \emph{dobro} de viagens do que seria eficiente --- alternativa (B). A sobreutiliza\c{c}\~ao t\'ipica do recurso comum.

\textbf{Distratores.} A alternativa (A), $N^{\text{AL}}=30,\,N^{*}=60$, \emph{inverte} os pap\'eis de m\'edio e marginal: aplica o marginal ao acesso livre e o m\'edio ao \'otimo --- exatamente o erro conceitual que a quest\~ao isola, pois o motorista individual captura o m\'edio, n\~ao o marginal. A alternativa (C), $N^{\text{AL}}=N^{*}=60$, sup\~oe que acesso livre e \'otimo coincidem, ignorando a externalidade de congestionamento; s\'o valeria com benef\'icio marginal igual ao m\'edio (curva plana), o que aqui \'e falso. As demais erram a \'algebra das igualdades.

\textbf{Intui\c{c}\~ao econ\^omica:} cada motorista que entra na zona central leva para casa o \emph{retorno m\'edio} da via, mas impõe a todos os demais o custo do congestionamento extra --- um custo marginal que n\~ao paga. Como o privado v\^e o m\'edio (acima do marginal enquanto $B$ \'e c\^oncavo), a entrada vai al\'em do ponto eficiente: $60$ em vez de $30$ mil viagens. O ped\'agio de Manhattan \'e justamente a cunha que faz o motorista enxergar o custo que impõe aos demais, recuando o tr\'afego ao n\'ivel eficiente.
\end{solucao}

\textbf{c)} (2 pontos, V ou F) A lama de rejeitos do rompimento da barragem de Funda\~ao (Samarco) percorreu o Rio Doce de Minas Gerais ao Esp\'irito Santo; em outubro de 2024 firmou-se um acordo de repactua\c{c}\~ao de cerca de R\$~170 bilh\~oes. Trata-se de uma externalidade negativa \emph{interestadual}: o dano gerado em um estado recai tamb\'em sobre o outro, rio abaixo. \textbf{Afirma\c{c}\~ao:} como o dano \'e interestadual, basta que o estado onde se localiza a fonte poluidora institua, agindo \emph{isoladamente}, um imposto pigouviano sobre a atividade em seu territ\'orio para alcan\c{c}ar o \'otimo de Pareto, pois um ente que tributa a fonte do dano necessariamente internaliza a externalidade que ela gera.

\gabaritoitem{F}
\begin{solucao}
{\footnotesize\textit{Paralelo de: Simulado 3, Bloco 3, VF2 (externalidade interestadual / queimadas de 2024 — agora ancorada em incentivo jurisdicional, não em cobertura parcial).}}\par\vspace{0.2em}
\textbf{Passo 1 --- o que se afirma.} Que um imposto pigouviano fixado por um \emph{\'unico} estado --- aquele que abriga a fonte poluidora --- agindo isoladamente, basta para alcan\c{c}ar o \'otimo de Pareto, sob o argumento de que tributar a fonte j\'a internaliza a externalidade.

\textbf{Passo 2 --- o que um ente isolado de fato internaliza.} A corre\c{c}\~ao pigouviana exige que a fonte enfrente o custo social marginal \emph{total} que gera --- aqui, a soma do dano que recai sobre o pr\'oprio estado-fonte (Minas Gerais) e do dano que desce o rio at\'e o Esp\'irito Santo. Mas um estado que legisla isoladamente tem incentivo a calibrar a al\'iquota apenas pelo dano que sente \emph{dentro} de sua jurisdi\c{c}\~ao: o preju\'izo do vizinho rio abaixo n\~ao entra em seu c\'alculo pol\'itico nem aparece em seu or\c{c}amento. Ele subtributa do ponto de vista social, fixando $\tau$ igual ao dano \emph{dom\'estico}, menor que o dano social total.

\textbf{Passo 3 --- por que \'e falsa.} ``Tributar a fonte do dano'' n\~ao implica ``internalizar o dano total''. O imposto s\'o conduz ao \'otimo se a al\'iquota igualar o dano marginal social \emph{somado sobre os dois estados}; um ente isolado, otimizando apenas o pr\'oprio bem-estar, n\~ao tem motivo para incluir o dano alheio. A atividade poluidora termina acima do n\'ivel eficiente, porque o pre\c{c}o-sombra que a fonte enfrenta cobre s\'o parte do estrago. Corrigir uma externalidade que cruza fronteiras exige coordena\c{c}\~ao interestadual, compet\^encia federal ou um acordo que internalize o dano \emph{ao longo de todo o curso do rio} --- n\~ao decorre de uma a\c{c}\~ao unilateral. A afirma\c{c}\~ao \'e, portanto, \textbf{falsa}.

\textbf{Passo 4 --- armadilha.} O erro t\'ipico \'e supor que, por estar a fonte concentrada num \'unico estado, basta esse estado tribut\'a-la para resolver o problema. A localiza\c{c}\~ao da fonte n\~ao muda o ponto: o que importa para a efici\^encia \'e que a al\'iquota cubra o dano \emph{total}, e o ente isolado n\~ao tem incentivo a fix\'a-la nesse n\'ivel --- internaliza apenas o que o atinge. \'E justamente o que distingue uma a\c{c}\~ao unilateral de um pacto federativo como o de R\$~170 bilh\~oes, negociado com a Uni\~ao e os dois estados.

\textbf{Intui\c{c}\~ao econ\^omica:} de nada adianta tributar a fonte se quem fixa a al\'iquota s\'o sente parte do preju\'izo. A lama desce o Rio Doce sem respeitar a divisa entre os estados, mas a vontade de tributar de cada ente p\'ara na fronteira do pr\'oprio dano. Corrigir uma externalidade interestadual pede um instrumento cuja jurisdi\c{c}\~ao --- e cujos incentivos --- cubram o curso inteiro do rio; da\'i o papel da Uni\~ao e de pactos federativos, e n\~ao de uma a\c{c}\~ao isolada do estado-fonte, por mais bem-intencionada que seja.
\end{solucao}

\textbf{d)} (2 pontos, V ou F) A previs\~ao meteorol\'ogica produzida e divulgada pelo INMET --- alertas de tempo severo, prognose de chuvas --- \'e um bem p\'ublico puro: n\~ao-rival e n\~ao-exclu\'ivel. Considere dois agricultores, $1$ e $2$, com disposi\c{c}\~oes marginais a pagar pelo n\'ivel $G$ de cobertura do servi\c{c}o dadas por $\text{DMP}^1(G)=10-G$ e $\text{DMP}^2(G)=6-G$ (em R\$, para $G\le 6$). \textbf{Afirma\c{c}\~ao:} como o servi\c{c}o \'e n\~ao-rival, a altura da curva de demanda \emph{agregada} num n\'ivel $G$ fixo \'e a \emph{soma} das duas DMP naquele $G$ --- por exemplo, em $G=2$ vale $R\$~12$ ($8+4$) ---, e n\~ao a soma das quantidades a um pre\c{c}o fixo, como seria num bem privado.

\gabaritoitem{V}
\begin{solucao}
{\footnotesize\textit{Paralelo de: Simulado 3, Bloco 3, VF1 (soma vertical vs.}}\par\vspace{0.2em}
\textbf{Passo 1 --- o que se afirma.} Que, por n\~ao-rivalidade, a curva de demanda agregada por bem p\'ublico se obt\'em \emph{empilhando} (somando verticalmente) as DMP individuais a cada $G$ fixo, com o valor num\'erico em $G=2$ igual a $R\$~12$.

\textbf{Passo 2 --- verificar a conta.} Em $G=2$: $\text{DMP}^1(2)=10-2=8$ e $\text{DMP}^2(2)=6-2=4$. Como o \emph{mesmo} n\'ivel $G=2$ \'e consumido simultaneamente pelos dois agricultores (n\~ao-rivalidade), o benef\'icio social marginal dessa unidade \'e o quanto \emph{ambos} pagariam por ela: $8+4=12$. A altura agregada em $G=2$ \'e, portanto, $R\$~12$, exatamente como a afirma\c{c}\~ao indica.

\textbf{Passo 3 --- por que vertical, e n\~ao horizontal.} Num bem privado, cada unidade vendida vai para \emph{um} consumidor: a um pre\c{c}o fixo somam-se as \emph{quantidades} que cada um quer (soma horizontal). Num bem p\'ublico puro n\~ao faz sentido somar quantidades, pois a quantidade $G$ \'e comum a todos; o que se soma s\~ao os \emph{pre\c{c}os} (DMP) a cada $G$ fixo --- soma vertical. Geometricamente, empilham-se as alturas no eixo dos pre\c{c}os a cada ponto do eixo das quantidades. \'E a leitura da condi\c{c}\~ao de Samuelson, $\sum_i\text{TMS}^i=\text{TMT}$: a demanda agregada vertical $\text{DMP}^1+\text{DMP}^2=16-2G$ cruza o custo marginal no n\'ivel eficiente.

\textbf{Passo 4 --- conclus\~ao.} A afirma\c{c}\~ao descreve corretamente a soma vertical (altura agregada $=$ soma das DMP a $G$ fixo) e confirma o valor $R\$~12$ em $G=2$, contrastando-a com a soma horizontal do bem privado. Logo \'e \textbf{verdadeira}.

\textbf{Intui\c{c}\~ao econ\^omica:} num bem privado, se dois agricultores querem fertilizante, o mercado precisa de dois lotes --- somam-se as quantidades. Num bem p\'ublico, se os dois se beneficiam do \emph{mesmo} alerta de geada do INMET, o alerta \'e um s\'o, mas o valor social dele \'e a soma do que cada um lhe atribui --- somam-se as valora\c{c}\~oes. Por isso a demanda agregada por bem p\'ublico empilha disposi\c{c}\~oes a pagar (vertical), e n\~ao quantidades (horizontal): a mesma unidade serve a todos simultaneamente, e \'e essa simultaneidade que justifica somar as DMP no eixo dos pre\c{c}os.
\end{solucao}

\quadroresposta{3}

\textbf{Questão 6.} (15 pontos) A Sabesp, privatizada em 2024, retomou com vigor o programa de despolui\c{c}\~ao do Rio Tiet\^e na regi\~ao metropolitana de S\~ao Paulo. Considere um trecho do rio cuja despolui\c{c}\~ao beneficia diretamente duas empresas ribeirinhas, $A$ (uma marina) e $B$ (um clube n\'autico), que captam recrea\c{c}\~ao e valor imobili\'ario da \'agua limpa. A qualidade da \'agua $G$ \'e um bem p\'ublico local: n\~ao-rival (a \'agua limpa que serve a marina serve igualmente o clube) e n\~ao-exclu\'ivel dentro do trecho. Cada empresa tem prefer\^encias quase-lineares sobre um bem privado $x^i$ (em R\$~mil) e o n\'ivel de despolui\c{c}\~ao $G$: suponha $u^A(x^A,G)=x^A+4\ln G$ e $u^B(x^B,G)=x^B+12\ln G$. Cada empresa disp\~oe de uma dota\c{c}\~ao inicial $\omega^i$ de bem privado, suficientemente grande para garantir solu\c{c}\~ao interior em $x^i$; como $x^i$ entra linearmente na utilidade, o n\'ivel de $\omega^i$ n\~ao afeta as condi\c{c}\~oes de otimalidade em $G$. Cada unidade de $G$ custa uma unidade de bem privado, de modo que $\text{TMT}_{G,x}=1$. Trate $G$ como cont\'inuo e estritamente positivo.

\textbf{a)} (5 pontos) Escreva a taxa marginal de substitui\c{c}\~ao $\text{TMS}^i_{G,x}$ de cada empresa e use a condi\c{c}\~ao de Samuelson para encontrar o n\'ivel socialmente \'otimo $G^*$ de despolui\c{c}\~ao.
\resolucaobox{7.2cm}
\begin{solucao}
{\footnotesize\textit{Paralelo de: Simulado 1, Bloco 3, Aberta (bem público ln G: Samuelson $\to$ Lindahl $\to$ provisão voluntária/carona; condomínio de segurança).}}\par\vspace{0.2em}
\textbf{Passo 1 --- dados.} $u^A=x^A+4\ln G$, $u^B=x^B+12\ln G$, $\text{TMT}_{G,x}=1$. Pede-se $\text{TMS}^i$ e $G^*$.

\textbf{Passo 2 --- TMS.} Para prefer\^encias quase-lineares, $\text{TMS}^i_{G,x}=\dfrac{\partial u^i/\partial G}{\partial u^i/\partial x^i}$. Como $\partial u^i/\partial x^i=1$ e $\partial u^A/\partial G=4/G$, $\partial u^B/\partial G=12/G$:
\[
\text{TMS}^A_{G,x}=\frac{4}{G},\qquad \text{TMS}^B_{G,x}=\frac{12}{G}.
\]

\textbf{Passo 3 --- Samuelson.} A provis\~ao eficiente de bem p\'ublico iguala a \emph{soma} das TMS \`a TMT:
\[
\text{TMS}^A+\text{TMS}^B=\text{TMT}\;\Longrightarrow\;\frac{4}{G}+\frac{12}{G}=1\;\Longrightarrow\;\frac{16}{G}=1\;\Longrightarrow\; G^{*}=16.
\]

\textbf{Passo 4 --- verifica\c{c}\~ao.} $G^{*}=16>0$, coerente com $\ln G$ definido. Dimens\~ao: TMS \'e adimensional (raz\~ao de utilidades marginais) e iguala TMT, tamb\'em adimensional. \textbf{(Passo 5 --- peer-review)} a soma \emph{vertical} das disposi\c{c}\~oes marginais a pagar --- e n\~ao a horizontal, como em bem privado --- \'e o ponto que o distrator t\'ipico erra; aqui est\'a aplicada corretamente, e a raz\~ao de valora\c{c}\~oes $4{:}12=1{:}3$ produz $G^{*}=16$ sem coincid\^encias num\'ericas esp\'urias.

\textbf{Intui\c{c}\~ao econ\^omica:} como a \'agua limpa \'e n\~ao-rival, a \emph{mesma} unidade de $G$ \'e desfrutada pelas duas empresas ao mesmo tempo. O valor social de uma unidade marginal \'e a soma do que cada uma pagaria por ela; igualar essa soma ao custo marginal ($\text{TMT}=1$) d\'a a regra de Samuelson e o n\'ivel $G^{*}=16$.
\rubricalabel
\begin{itemize}\item Obt\'em corretamente $\text{TMS}^A=4/G$ e $\text{TMS}^B=12/G$ (utilidade marginal de $G$ sobre a de $x$): 2 pts (1 pt cada).\item Enuncia/aplica a condi\c{c}\~ao de Samuelson como \emph{soma} das TMS igual \`a TMT (n\~ao a m\'edia, n\~ao individual): 2 pts.\item Resolve e chega a $G^{*}=16$: 1 pt.\item Erro aritm\'etico isolado com m\'etodo correto (ex.: soma certa mas conta errada de $16/G=1$): penalidade leve de 1 pt no total do sub-item.\item Usar soma horizontal (estilo bem privado) ou igualar TMS individual \`a TMT: m\'aximo 1 pt no sub-item.\item \emph{N\~ao-cumulatividade:} tetos por erro conceitual (ex.: soma horizontal $\to$ m\'ax 1 pt) prevalecem sobre penalidades aritm\'eticas leves --- n\~ao se somam.\end{itemize}
\end{solucao}

\textbf{b)} (5 pontos) Para descentralizar $G^{*}$ via um equil\'ibrio de Lindahl, cada empresa paga uma fra\c{c}\~ao personalizada $\tau^i$ do custo de cada unidade de $G$. Determine os pre\c{c}os de Lindahl $(\tau^A,\tau^B)$ e verifique que eles ``somam'' a TMT. Comente em uma frase por que esses pre\c{c}os diferem entre as empresas.
\resolucaobox{7.2cm}
\begin{solucao}
\textbf{Passo 1 --- plano.} No equil\'ibrio de Lindahl, cada empresa enfrenta seu pre\c{c}o personalizado $\tau^i$ por unidade e escolhe, voluntariamente, o \emph{mesmo} $G$ --- que deve coincidir com $G^{*}$. Isso exige $\tau^i=\text{TMS}^i$ avaliado em $G^{*}$, com $\sum_i\tau^i=\text{TMT}$.

\textbf{Passo 2 --- c\'alculo.} Avaliando as TMS em $G^{*}=16$:
\[
\tau^A=\text{TMS}^A_{G,x}(G^{*})=\frac{4}{16}=\frac{1}{4},\qquad \tau^B=\text{TMS}^B_{G,x}(G^{*})=\frac{12}{16}=\frac{3}{4}.
\]

\textbf{Passo 3 --- verifica\c{c}\~ao.} $\tau^A+\tau^B=\tfrac{1}{4}+\tfrac{3}{4}=1=\text{TMT}$. As fra\c{c}\~oes somam o custo unit\'ario total, como exige Lindahl. Confer\^encia de otimalidade: a empresa $A$, sob pre\c{c}o $1/4$, maximiza $x^A+4\ln G$ com $x^A=\omega^A-\tfrac{1}{4}G$, CPO $4/G=1/4\Rightarrow G=16$; e $B$, sob $3/4$, CPO $12/G=3/4\Rightarrow G=16$. Ambas concordam em $G=16$. \textbf{(Passo 5 --- peer-review)} risco residual: o aluno pode confundir $\tau^i$ com o \emph{valor pago total} $\tau^i G$; o enunciado pede o \emph{pre\c{c}o} (fra\c{c}\~ao por unidade), evitando ambiguidade.

\textbf{Intui\c{c}\~ao econ\^omica:} Lindahl cobra de cada uma \emph{na propor\c{c}\~ao do quanto valoriza} a \'agua limpa. O clube $B$, que valoriza a despolui\c{c}\~ao o triplo de $A$ ($12$ contra $4$), paga o triplo por unidade ($3/4$ contra $1/4$). Com pre\c{c}os assim personalizados, cada uma, agindo egoisticamente, escolhe exatamente o n\'ivel eficiente --- os pre\c{c}os diferem porque as disposi\c{c}\~oes marginais a pagar diferem.
\rubricalabel
\begin{itemize}\item Reconhece que $\tau^i=\text{TMS}^i$ avaliado em $G^{*}=16$: 2 pts.\item Calcula $\tau^A=1/4$ e $\tau^B=3/4$: 2 pts (1 pt cada).\item Verifica $\tau^A+\tau^B=1=\text{TMT}$ e justifica em uma frase a diferen\c{c}a ($B$ valoriza mais, logo paga mais por unidade): 1 pt.\item Avaliar a TMS no $G$ errado (ex.: $G=12$) mas com m\'etodo correto: penalidade leve de 1 pt.\item Apresentar $\tau^i$ como valor pago total $\tau^i G$ em vez de pre\c{c}o por unidade: m\'aximo 3 pts no sub-item.\item \emph{N\~ao-cumulatividade:} tetos por erro conceitual (ex.: $\tau^i$ como valor total $\to$ m\'ax 3 pts) prevalecem sobre penalidades aritm\'eticas leves --- n\~ao se somam.\end{itemize}
\end{solucao}

\textbf{c)} (5 pontos) Suponha agora que n\~ao haja mecanismo de Lindahl: cada empresa decide \emph{voluntariamente} quanto contribuir para a despolui\c{c}\~ao, $g^i\ge 0$, com $G=g^A+g^B$, tomando a contribui\c{c}\~ao da outra como dada (provis\~ao volunt\'aria / problema do carona). Determine o n\'ivel total provido $G^N$, quem contribui e quem ``pega carona'', compare $G^N$ com $G^{*}$ e explique economicamente o resultado.
\resolucaobox{7.2cm}
\begin{solucao}
\textbf{Passo 1 --- montagem.} Cada empresa resolve $\max_{g^i\ge 0}\;u^i=(\omega^i-g^i)+\theta^i\ln(g^i+g^{-i})$, com $\theta^A=4$, $\theta^B=12$, tomando $g^{-i}$ como dado (melhor-resposta). Como $\omega^i$ \'e constante, ela some na derivada; a CPO interior \'e $\dfrac{\theta^i}{g^i+g^{-i}}=1$, ou seja, cada empresa \emph{deseja} um total $G$ igual a $\theta^i$.

\textbf{Passo 2 --- rea\c{c}\~ao \'otima.} A demanda individual por \emph{total} \'e $G=\theta^A=4$ para $A$ e $G=\theta^B=12$ para $B$. Como as duas n\~ao podem valer simultaneamente, vale a maior: a empresa de maior valora\c{c}\~ao ($B$, $\theta^B=12$) provê at\'e seu \'otimo individual, e a de menor ($A$) encontra o total j\'a acima do que desejaria e contribui zero.

\textbf{Passo 3 --- equil\'ibrio.} Verifica-se: se $g^B=12$ e $g^A=0$, ent\~ao $G=12$. Para $A$: $\theta^A/G=4/12=1/3<1$, logo a utilidade marginal de contribuir ($4/G$) \'e menor que o custo ($1$) --- $A$ n\~ao contribui, \'otimo $g^A=0$ (solu\c{c}\~ao de canto). Para $B$: $\theta^B/G=12/12=1$, CPO satisfeita, $B$ n\~ao quer aumentar nem reduzir. Assim
\[
g^A=0,\quad g^B=12,\quad G^N=12.
\]
Comparando, $G^N=12<G^{*}=16$; a raz\~ao \'e $G^N/G^{*}=12/16=3/4$. Subprovis\~ao.

\textbf{Passo 4 --- verifica\c{c}\~ao e peer-review.} Casos-limite: a empresa de menor valora\c{c}\~ao \'e quem pega carona, consistente com a literatura de provis\~ao volunt\'aria assim\'etrica (Bergstrom, Blume \& Varian, 1986). \textbf{(Passo 5)} risco residual: a forma logar\'itmica garante CPO interior bem definida e canto limpo em $g^A=0$ --- com $\sqrt{G}$ surgiriam coincid\^encias num\'ericas esp\'urias entre m\'edia das valora\c{c}\~oes e Samuelson; o curso fixa $\ln G$ justamente para evit\'a-las. Confere.

\textbf{Intui\c{c}\~ao econ\^omica:} sob contribui\c{c}\~ao volunt\'aria, cada empresa s\'o conta o pr\'oprio benef\'icio marginal --- ignora o que sua contribui\c{c}\~ao vale para a outra. A de maior valora\c{c}\~ao acaba ``segurando'' a despolui\c{c}\~ao sozinha, e a de menor pega carona ($g^A=0$). O total $G^N=12$ fica abaixo do eficiente $G^{*}=16$ porque a soma das TMS no n\'ivel provido, $16/12=4/3>1$, ainda excede o custo: a sociedade valoriza mais \'agua limpa do que o mercado descentralizado entrega. \'E exatamente a falha que Lindahl, ou um rateio compuls\'orio, corrige.
\rubricalabel
\begin{itemize}\item Monta o problema de provis\~ao volunt\'aria com $G=g^A+g^B$ e deriva a CPO $\theta^i/G=1$ (cada empresa deseja $G=\theta^i$): 2 pts.\item Identifica que apenas $B$ contribui ($g^B=12$) e $A$ pega carona ($g^A=0$), justificando o canto via $\theta^A/G<1$: 1{,}5 pt.\item Conclui $G^N=12$, compara com $G^{*}=16$ e aponta subprovis\~ao (raz\~ao $3/4$): 1 pt.\item Explica\c{c}\~ao econ\^omica do free-rider (cada um ignora o benef\'icio externo da pr\'opria contribui\c{c}\~ao): 0{,}5 pt.\item Erro aritm\'etico isolado com m\'etodo de provis\~ao volunt\'aria correto: penalidade leve de 1 pt.\item Tratar como solu\c{c}\~ao interior com ambas contribuindo positivamente (ignorando o canto) e obter $G^N\neq 12$: m\'aximo 2{,}5 pts no sub-item.\item \emph{N\~ao-cumulatividade:} tetos por erro conceitual (ex.: ignorar o canto $\to$ m\'ax 2{,}5 pts) prevalecem sobre penalidades aritm\'eticas leves --- n\~ao se somam.\end{itemize}
\end{solucao}



\section*{Bloco 4 --- Assimetria Informacional e Sinalização \hfill \normalsize (25 pontos)}

\textbf{Questão 7.} (10 pontos) \textbf{Itens objetivos.} Nos itens \textbf{a)} e \textbf{b)}, assinale a \emph{única} alternativa correta (3 pontos cada; não há penalização por erro). Nos itens \textbf{c)} e \textbf{d)}, responda \textbf{Verdadeiro ou Falso} (2 pontos cada; vale a regra de anulação descrita na capa). Registre suas respostas no quadro-resposta ao final do bloco.

\textbf{a)} (3 pontos) O mercado brasileiro de \emph{recommerce} (revenda de smartphones recondicionados) cresceu com plataformas que intermedeiam a venda de aparelhos seminovos: cada vendedor conhece a qualidade $q$ do pr\'oprio aparelho, mas o comprador n\~ao. Suponha que a qualidade seja distribu\'ida uniformemente, $q \sim U[0,30]$ (em centenas de reais), que cada vendedor s\'o aceite vender por um pre\c{c}o que cubra ao menos $\tfrac{3}{4}$ da qualidade (reserva $\tfrac{3}{4}q$), e que um comprador neutro ao risco valore um aparelho de qualidade $q$ em $k\,q$, com $k>0$. A plataforma posta um pre\c{c}o \'unico $p>0$ e o comprador adquire de quem aceitar. Para qual condi\c{c}\~ao sobre $k$ existe pre\c{c}o positivo com transa\c{c}\~oes em equil\'ibrio (o mercado n\~ao colapsa)?

\begin{enumerate}[label=\Alph*., leftmargin=3.2em, itemsep=0.15em, topsep=0.2em]
  \item $k \ge \tfrac{3}{4}$
  \item $k \ge 2$
  \item $k \ge \tfrac{3}{2}$
  \item $k \ge 3$
  \item $k > 0$ (qualquer pr\^emio positivo basta)
\end{enumerate}
\gabaritoitem{C}
\begin{solucao}
{\footnotesize\textit{Paralelo de: Simulado 1, Bloco 4, ME1 (Akerlof/lemons, Webmotors): média condicional de uniforme truncada e unraveling, condição sobre o prêmio k para o mercado não colapsar.}}\par\vspace{0.2em}
\textbf{Passo 1 (quem aceita vender).} Ao pre\c{c}o $p$, vende apenas quem tem reserva $\tfrac{3}{4}q \le p$, isto \'e, os aparelhos com $q \le \tfrac{4}{3}p$. Para $p$ interior $\big(\tfrac{4}{3}p \le 30\big)$, o conjunto que aparece para vender \'e $q \in \big[0,\tfrac{4}{3}p\big]$, e pela uniformidade a qualidade m\'edia \emph{condicional} aos que aceitam \'e
\[
E\!\left[q \,\middle|\, q \le \tfrac{4}{3}p\right] = \frac{0 + \tfrac{4}{3}p}{2} = \frac{2p}{3}.
\]

\textbf{Passo 2 (disposi\c{c}\~ao a pagar do comprador).} O comprador racional antecipa a sele\c{c}\~ao adversa: o que aparece para venda \'e pior que a popula\c{c}\~ao total. Seu valor esperado \'e $k\,E\!\left[q \mid q \le \tfrac{4}{3}p\right] = k\,\dfrac{2p}{3}$. Ele s\'o paga $p$ se
\[
k\,\frac{2p}{3} \;\ge\; p \quad\Longleftrightarrow\quad k \ge \frac{3}{2}\quad(\text{para } p>0).
\]

\textbf{Passo 3 (conclus\~ao).} Se $k < \tfrac{3}{2}$, a desigualdade falha para todo $p>0$ e o \'unico equil\'ibrio \'e $p^{*}=0$ (colapso por \emph{unraveling}). Se $k \ge \tfrac{3}{2}$, h\'a pre\c{c}o positivo sustent\'avel. A condi\c{c}\~ao \'e $k \ge \tfrac{3}{2}$.

\textbf{Por que os distratores mais tentadores falham.} A alternativa $k \ge 2$ \'e o erro de quem \emph{esquece o fator de reserva} $\tfrac{3}{4}$: trata a reserva como $q$ cheio, recupera a m\'edia condicional $p/2$ e obt\'em $k \ge 2$ \textemdash{} o resultado do caso-livro, mas n\~ao deste enunciado. A alternativa $k \ge \tfrac{3}{4}$ \'e o erro de quem \emph{ignora o truncamento por completo} e l\^e o coeficiente de reserva $\tfrac{3}{4}$ do enunciado como se fosse o pr\'oprio limiar: ao comparar $k\,q$ com a reserva $\tfrac{3}{4}q$ sem aplicar o desconto de sele\c{c}\~ao, o fator $q$ cancela e sobra $k \ge \tfrac{3}{4}$ \textemdash{} um n\'umero que aparece s\'o porque \'e o coeficiente do enunciado, n\~ao porque resolva o problema. A alternativa $k \ge 3$ confunde a m\'edia da uniforme em $\big[0,\tfrac{4}{3}p\big]$, tomando-a como $\tfrac{1}{4}$ do extremo em vez de $\tfrac{1}{2}$.

\textbf{Intui\c{c}\~ao econ\^omica:} qualquer pre\c{c}o positivo seleciona para baixo a qualidade ofertada, e o comprador, sabendo disso, desconta. A m\'edia condicional dos que vendem \'e fra\c{c}\~ao $\tfrac{2}{3}$ do pre\c{c}o \textemdash{} o vi\'es de truncamento $\tfrac{1}{2}$ combinado ao fator de reserva $\tfrac{4}{3}$, pois $q \le \tfrac{4}{3}p$. O mercado s\'o sobrevive quando o pr\^emio $k$ compensa esse desconto, isto \'e, $k \ge \dfrac{1}{2/3} = \tfrac{3}{2}$ \textemdash{} note que o limiar \'e o \emph{rec\'iproco} da fra\c{c}\~ao $\tfrac{2}{3}$, n\~ao o seu valor. Reservas mais baixas (vendedores menos exigentes) suavizam o problema: aqui o limiar cai de $2$ para $\tfrac{3}{2}$ justamente porque a reserva \'e fra\c{c}\~ao da qualidade, n\~ao a qualidade inteira.
\end{solucao}

\textbf{b)} (3 pontos) A Casas Bahia vende garantia estendida (subscrita por seguradora parceira) sobre eletr\^onicos da Samsung. Cada fabricante conhece a robustez do pr\'oprio aparelho; o consumidor, n\~ao. H\'a dois tipos, com probabilidade de falha dentro do per\'iodo de garantia dada por $p_R = 0{,}2$ (robusto) e $p_F = 0{,}6$ (fr\'agil). Um aparelho percebido como ``robusto'' vende com pr\^emio $\Delta = R\$~9$ (em centenas de reais) sobre um ``fr\'agil''. Oferecer a garantia compromete o fabricante a pagar $B$ (em centenas de reais) a cada falha, com custo esperado $p_i\,B$ para o tipo $i$. Para a garantia funcionar como sinal separador (o robusto oferece, o fr\'agil n\~ao imita), qual o menor $B$?

\begin{enumerate}[label=\Alph*., leftmargin=3.2em, itemsep=0.15em, topsep=0.2em]
  \item $B = 9$
  \item $B = 45$
  \item $B = 22{,}5$
  \item Qualquer $B>0$ separa, pois a garantia sempre custa mais ao tipo fr\'agil
  \item $B = 15$
\end{enumerate}
\gabaritoitem{E}
\begin{solucao}
{\footnotesize\textit{Paralelo de: Simulado 3, Bloco 4, ME1 (garantia como sinal, BYD/bateria): menor sinal B=$\Delta$/p\_L que separa, com checagem do teto $\Delta$/p\_H.}}\par\vspace{0.2em}
\textbf{Passo 1 (por que a garantia sinaliza).} O custo esperado de oferecer a garantia \'e $p_i B$, maior para quem tem aparelho pior: $p_F B > p_R B$. O sinal separa se o fabricante fr\'agil \emph{n\~ao} achar vantajoso imit\'a-lo.

\textbf{Passo 2 (deriva\c{c}\~ao do piso).} Imitar significa oferecer a garantia $B$ e capturar o pr\^emio $\Delta$ ao custo esperado $p_F B$. O fr\'agil imita se e somente se $\Delta > p_F B$. Logo a n\~ao-imita\c{c}\~ao exige
\[
p_F\,B \;\ge\; \Delta \quad\Longleftrightarrow\quad B \;\ge\; \frac{\Delta}{p_F} = \frac{9}{0{,}6} = 15.
\]

\textbf{Passo 3 (o robusto aceita).} O robusto oferece a garantia desde que seu custo esperado n\~ao supere o pr\^emio: $p_R B \le \Delta \Leftrightarrow B \le \Delta/p_R = 9/0{,}2 = 45$. O intervalo de sinais cr\'iveis \'e $[15,\,45]$, n\~ao-vazio porque $p_F > p_R$. O menor sinal separador \'e $B = 15$.

\textbf{Por que os distratores mais tentadores falham.} A alternativa $B = 45$ usa $\Delta/p_R$, que \'e o \emph{teto} (at\'e onde o robusto ainda tolera oferecer), n\~ao o piso que dissuade o fr\'agil; \'e o erro de olhar a restri\c{c}\~ao do tipo errado. A alternativa $B = 9$ confunde o pr\^emio $\Delta$ com o pr\'oprio sinal $B$, esquecendo de dividir pela probabilidade de acionamento $p_F$ \textemdash{} ignora que a garantia s\'o \'e acionada com probabilidade $0{,}6$, de modo que para impor custo $9$ ao fr\'agil \'e preciso prometer pagar mais que $9$ por falha. (A alternativa $B = 22{,}5 = \Delta/\bar p$, com $\bar p = 0{,}4$ a m\'edia das probabilidades, usa o tipo m\'edio em vez do fr\'agil, que \'e quem de fato decide imitar.)

\textbf{Intui\c{c}\~ao econ\^omica:} a garantia \'e cr\'edivel como sinal porque d\'oi mais em quem tem o produto pior. O custo diferencial nasce da diferen\c{c}a de probabilidade de acionamento, e o sinal precisa ser caro o bastante para que o custo esperado do fr\'agil supere o pr\^emio de reputa\c{c}\~ao. Abaixo de $\Delta/p_F$, a garantia \'e barata demais at\'e para o fr\'agil, e ele imita.
\end{solucao}

\textbf{c)} (2 pontos, V ou F) O Sistema Nacional de Transplantes (SNT) aloca \'org\~aos a receptores combinando compatibilidade e prioridade cl\'inica \textemdash{} uma estrutura de \emph{matching} \`a la Roth. Considere dois receptores, Helena e Rafael, e dois rins dispon\'iveis, Rim~1 e Rim~2 (um para cada receptor). As prefer\^encias dos receptores s\~ao: Helena e Rafael \emph{ambos} preferem o Rim~1 ao Rim~2 (melhor compatibilidade esperada). As prioridades dos \'org\~aos s\~ao: o Rim~1 prioriza Helena $\succeq$ Rafael; o Rim~2 prioriza Rafael $\succeq$ Helena. Foi proposta, fora do algoritmo, a aloca\c{c}\~ao $\mu$: Helena $\to$ Rim~2 e Rafael $\to$ Rim~1. Afirma-se: \emph{o par (Helena, Rim~1) \'e um par bloqueante, de modo que $\mu$ n\~ao \'e est\'avel}.

\gabaritoitem{V}
\begin{solucao}
{\footnotesize\textit{Paralelo de: Simulado 2, Bloco 4, VF1 (estabilidade/par bloqueante, SiSU/Gale-Shapley): testar a DEFINIÇÃO de estabilidade via existência de um par bloqueante concreto.}}\par\vspace{0.2em}
\textbf{Passo 1 (crit\'erio de estabilidade).} Um \emph{matching} \'e est\'avel se n\~ao existe par bloqueante: um par receptor--\'org\~ao que \emph{n\~ao} est\'a casado entre si, mas em que ambos prefeririam estar um com o outro do que com seu par atual em $\mu$. O ju\'izo pedido \'e sobre a \emph{defini\c{c}\~ao} de estabilidade, n\~ao sobre qual lado o algoritmo favorece.

\textbf{Passo 2 (testar o par (Helena, Rim~1)).} Em $\mu$, Helena est\'a com o Rim~2 e o Rim~1 est\'a com Rafael.
\begin{itemize}
\item \emph{Helena prefere o Rim~1?} Sim \textemdash{} sua prefer\^encia \'e Rim~1 $\succeq$ Rim~2 (estrita), e em $\mu$ ela tem o Rim~2.
\item \emph{O Rim~1 prefere Helena?} Sim \textemdash{} o Rim~1 prioriza Helena $\succeq$ Rafael, e em $\mu$ ele est\'a com Rafael.
\end{itemize}
Logo Helena e o Rim~1 prefeririam, ambos, formar par entre si: $(\text{Helena},\text{Rim~1})$ \'e par bloqueante.

\textbf{Passo 3 (conclus\~ao sobre $\mu$).} Como existe par bloqueante, $\mu$ \textbf{n\~ao \'e est\'avel}. A afirma\c{c}\~ao est\'a correta. (Como contraste, o \emph{matching} est\'avel seria Helena $\to$ Rim~1 e Rafael $\to$ Rim~2: a\'i cada \'org\~ao fica com seu receptor priorit\'ario e nenhum par bloqueia.)

\textbf{Intui\c{c}\~ao econ\^omica:} a estabilidade \'e a condi\c{c}\~ao m\'inima para uma aloca\c{c}\~ao ``sobreviver'' a recursos e contesta\c{c}\~oes. Se Helena, com prioridade suficiente, quer o Rim~1, e o Rim~1 a prioriza sobre quem ocupa a vaga, ambos t\^em incentivo a furar a aloca\c{c}\~ao proposta. Uma aloca\c{c}\~ao assim n\~ao se sustenta \textemdash{} \'e exatamente esse tipo de inconsist\^encia que a aceita\c{c}\~ao diferida elimina ao produzir um \emph{matching} est\'avel.
\end{solucao}

\textbf{d)} (2 pontos, V ou F) O Cadastro Nacional de Ado\c{c}\~ao (SNA), do Conselho Nacional de Justi\c{c}a, organiza a aproxima\c{c}\~ao entre pretendentes habilitados e crian\c{c}as/adolescentes dispon\'iveis, e pode ser modelado como um \emph{matching} por aceita\c{c}\~ao diferida (Gale--Shapley). Suponha uma rodada em que os \emph{pretendentes habilitados} sejam o lado proponente (fazem propostas; o outro lado aceita ou recusa provisoriamente segundo prioridades). Afirma-se: \emph{com os pretendentes propondo, o algoritmo entrega o matching est\'avel que \'e o melhor poss\'ivel para o lado que recebe as propostas \textemdash{} cada agente desse lado fica com o melhor parceiro que tem em algum matching est\'avel}.

\gabaritoitem{F}
\begin{solucao}
{\footnotesize\textit{Paralelo de: Simulado 1, Bloco 4, VF2 (Gale-Shapley, ENARE/NRMP): teorema da otimalidade do lado proponente — o resultado é ótimo para quem propõe e péssimo para quem recebe, entre os estáveis.}}\par\vspace{0.2em}
\textbf{Passo 1 (teorema da otimalidade do proponente).} Na aceita\c{c}\~ao diferida, o resultado \'e est\'avel e \'otimo \emph{para o lado que prop\~oe}: cada agente do lado proponente obt\'em o melhor parceiro que consegue em \emph{qualquer} \emph{matching} est\'avel. Simultaneamente, esse mesmo resultado \'e o \emph{pior} (entre os est\'aveis) para o lado que recebe as propostas.

\textbf{Passo 2 (aplica\c{c}\~ao).} Se os \emph{pretendentes} prop\~oem, o \emph{matching} gerado \'e \emph{\'otimo para os pretendentes} e \emph{p\'essimo para o lado receptor} (cada agente que recebe fica com o pior parceiro que aparece em algum \emph{matching} est\'avel). A afirma\c{c}\~ao troca os pap\'eis: descreve a otimalidade do lado \emph{receptor}, que s\'o ocorreria se fosse \emph{ele} a propor. Logo \'e \textbf{falsa}.

\textbf{Passo 3 (qual lado favorecer).} A escolha do lado proponente n\~ao \'e neutra \textemdash{} \'e uma decis\~ao distributiva sobre qual lado do mercado a institui\c{c}\~ao quer favorecer. Inverter quem prop\~oe inverte a quem o resultado \'otimo \'e entregue.

\textbf{Intui\c{c}\~ao econ\^omica:} quem prop\~oe controla a din\^amica \textemdash{} percorre sua lista de prefer\^encias de cima para baixo e s\'o desce quando \'e rejeitado, garantindo o melhor parceiro alcan\c{c}\'avel. Quem recebe apenas acumula propostas e fica preso \`a melhor que chegou, que pode ser fraca. Por isso a otimalidade pertence sempre ao lado proponente, e atribu\'i-la ao receptor inverte o teorema.
\end{solucao}

\quadroresposta{4}

\textbf{Questão 8.} (15 pontos) Com a expans\~ao do seguro param\'etrico de granizo para vinhedos da Serra Ga\'ucha, considere um operador monopolista que oferece a um vinhedo um \emph{n\'ivel de cobertura} $q \ge 0$ (\'indice de prote\c{c}\~ao param\'etrica contratada), cobrando um pagamento total $T$. A cobertura gera ao vinhedo um valor bruto $\theta\,q$, e seu excedente \'e $\theta q - T$. O custo do operador para entregar $q$ \'e $C(q) = \tfrac{1}{2}q^{2}$. H\'a dois tipos de vinhedo quanto \`a exposi\c{c}\~ao ao granizo, $\theta \in \{\theta_L,\theta_H\}$ com $\theta_L = 8$ e $\theta_H = 10$, cada um com probabilidade $\tfrac{1}{2}$. \emph{Cada vinhedo conhece o pr\'oprio $\theta$; o operador, n\~ao.} Todo vinhedo tem excedente de reserva normalizado a zero (recusa o contrato se o excedente for negativo). O operador desenha o contrato para maximizar o lucro esperado. Suponha que esses par\^ametros sejam meramente ilustrativos.

\textbf{a)} (5 pontos) Suponha que o operador \emph{observe} o tipo $\theta$ de cada vinhedo (\emph{first-best}, com discrimina\c{c}\~ao perfeita). Determine os n\'iveis de cobertura $q_L^{FB}, q_H^{FB}$, os pagamentos $T_L^{FB}, T_H^{FB}$ que extraem todo o excedente, e o lucro esperado do operador. Explique por que, sob informa\c{c}\~ao perfeita, n\~ao h\'a ``renda informacional''.
\resolucaobox{7.2cm}
\begin{solucao}
{\footnotesize\textit{Paralelo de: Simulado 3, Bloco 4, aberta (screening monopólico FB vs SB, data center ByteDance/Casa dos Ventos).}}\par\vspace{0.2em}
\textbf{Passo 1 (ponto-chave).} Observando o tipo, o operador trata cada vinhedo como um problema separado: escolhe a cobertura eficiente (valor marginal igual a custo marginal) e cobra o pagamento que zera o excedente do cliente, capturando todo o ganho de com\'ercio.

\textbf{Passo 2 (deriva\c{c}\~ao).} Para o tipo $\theta$, o operador maximiza $T - \tfrac{1}{2}q^{2}$ sujeito \`a participa\c{c}\~ao $\theta q - T \ge 0$. A restri\c{c}\~ao (\emph{IR}) vincula, logo $T = \theta q$ e o lucro vira $\theta q - \tfrac{1}{2}q^{2}$. A condi\c{c}\~ao de primeira ordem d\'a
\[
\theta - q = 0 \;\Longrightarrow\; q^{FB} = \theta.
\]
Portanto
\[
q_L^{FB} = 8,\quad q_H^{FB} = 10,\qquad T_L^{FB} = \theta_L q_L^{FB} = 64,\quad T_H^{FB} = \theta_H q_H^{FB} = 100.
\]
Lucro por tipo $T_i^{FB} - \tfrac{1}{2}(q_i^{FB})^{2}$: $64 - 32 = 32$ no tipo $L$ e $100 - 50 = 50$ no tipo $H$. Lucro esperado
\[
\Pi^{FB} = \tfrac{1}{2}\cdot 32 + \tfrac{1}{2}\cdot 50 = 41.
\]

\textbf{Passo 3 (verifica\c{c}\~ao).} As coberturas igualam valor marginal $\theta$ a custo marginal $q$ \textemdash{} s\~ao eficientes. Cada vinhedo fica com excedente exatamente zero ($\theta_i q_i - T_i = 0$), de modo que o operador captura todo o excedente social $\theta_i q_i - \tfrac{1}{2}q_i^{2}$.

\textbf{Intui\c{c}\~ao econ\^omica:} com o tipo observ\'avel, o operador faz a cada vinhedo uma oferta ``pegar ou largar'' sob medida. N\~ao h\'a informa\c{c}\~ao privada a proteger, ent\~ao n\~ao \'e preciso ceder \emph{renda informacional} a ningu\'em: a discrimina\c{c}\~ao \'e perfeita e a aloca\c{c}\~ao, eficiente.
\rubricalabel
Total 5 pontos. (i) 2 pontos pelas coberturas eficientes via valor marginal igual a custo marginal ($q_i^{FB} = \theta_i$, logo $q_L = 8,\ q_H = 10$); 1 ponto parcial se s\'o uma estiver correta ou se a CPO estiver montada mas mal resolvida. (ii) 1{,}5 ponto pelos pagamentos que extraem todo o excedente ($T_i = \theta_i q_i \Rightarrow T_L = 64,\ T_H = 100$). (iii) 1 ponto pelo lucro esperado $\Pi^{FB} = \tfrac{1}{2}(32) + \tfrac{1}{2}(50) = 41$. (iv) 0{,}5 ponto pela justificativa de aus\^encia de renda informacional: com tipo observ\'avel, o cliente n\~ao det\'em informa\c{c}\~ao privada, ent\~ao a \emph{IR} de cada tipo pode vincular e o excedente \'e todo capturado. Erro aritm\'etico isolado com m\'etodo correto (CPO certa, conta de lucro errada) perde no m\'aximo 0{,}5 ponto.
\end{solucao}

\textbf{b)} (5 pontos) Agora o operador \emph{n\~ao observa} o tipo (\emph{second-best}) e oferece um \emph{menu} $\{(q_L,T_L),$ $(q_H,T_H)\}$, deixando cada vinhedo escolher. Escreva as restri\c{c}\~oes de participa\c{c}\~ao (\emph{IR}) e de compatibilidade de incentivos (\emph{IC}) relevantes, identifique quais vinculam no \'otimo, e resolva para obter $q_L, q_H, T_L, T_H$ e a renda informacional do tipo $H$. Comente o resultado de ``aus\^encia de distor\c{c}\~ao no topo''.
\resolucaobox{7.2cm}
\begin{solucao}
\textbf{Passo 1 (ponto-chave).} Sem observar o tipo, o operador n\~ao pode mais zerar o excedente dos dois. Para impedir que o vinhedo $H$ se passe por $L$, precisa deixar-lhe uma \emph{renda informacional}; e para reduzir essa renda, distorce \emph{para baixo} a cobertura do tipo $L$. Vinculam a \emph{IR} do tipo baixo e a \emph{IC} do tipo alto.

\textbf{Passo 2 (restri\c{c}\~oes).} Excedente do tipo $\theta_i$ ao escolher o pacote $i$: $U_i = \theta_i q_i - T_i$.
\[
\textbf{(IR}_L\textbf{)}\ \theta_L q_L - T_L \ge 0,\qquad \textbf{(IR}_H\textbf{)}\ \theta_H q_H - T_H \ge 0,
\]
\[
\textbf{(IC}_H\textbf{)}\ \theta_H q_H - T_H \ge \theta_H q_L - T_L,\qquad \textbf{(IC}_L\textbf{)}\ \theta_L q_L - T_L \ge \theta_L q_H - T_H.
\]
No \'otimo vinculam $\textbf{(IR}_L\textbf{)}$ e $\textbf{(IC}_H\textbf{)}$; as outras duas ficam folgadas. De $\textbf{(IR}_L\textbf{)}$: $T_L = \theta_L q_L$, logo $U_L = 0$. De $\textbf{(IC}_H\textbf{)}$ com igualdade,
\[
U_H = \theta_H q_H - T_H = \theta_H q_L - T_L = (\theta_H - \theta_L)\,q_L,
\]
a renda informacional do tipo $H$.

\textbf{Passo 3 (deriva\c{c}\~ao).} Substituindo $T_H = \theta_H q_H - (\theta_H - \theta_L)q_L$ e $T_L = \theta_L q_L$ no lucro esperado,
\[
\Pi = \tfrac{1}{2}\!\left[\theta_H q_H - (\theta_H-\theta_L)q_L - \tfrac{1}{2}q_H^{2}\right] + \tfrac{1}{2}\!\left[\theta_L q_L - \tfrac{1}{2}q_L^{2}\right].
\]
CPO em $q_H$: $\theta_H - q_H = 0 \Rightarrow q_H = \theta_H = 10$ (eficiente). CPO em $q_L$:
\[
\tfrac{1}{2}(\theta_L - q_L) - \tfrac{1}{2}(\theta_H - \theta_L) = 0 \;\Longrightarrow\; q_L = \theta_L - (\theta_H - \theta_L) = 8 - 2 = 6.
\]
Logo
\[
q_H = 10,\quad q_L = 6,\qquad T_L = \theta_L q_L = 48,\qquad T_H = \theta_H q_H - (\theta_H-\theta_L)q_L = 100 - 2\cdot 6 = 88,
\]
e a renda informacional do tipo $H$ \'e $U_H = (\theta_H - \theta_L)q_L = 2\cdot 6 = 12$.

\textbf{Passo 4 (verifica\c{c}\~ao).} $\textbf{(IC}_H\textbf{)}$: no pacote $H$, $U_H = 10\cdot 10 - 88 = 12$; imitando $L$ seria $10\cdot 6 - 48 = 12$ \textemdash{} igual, vincula. $\textbf{(IC}_L\textbf{)}$: tipo $L$ no pacote $L$ tem $0$; imitando $H$ teria $8\cdot 10 - 88 = -8 < 0$ \textemdash{} folgada. $\textbf{(IR}_H\textbf{)}$: $U_H = 12 > 0$ \textemdash{} folgada. Coerente.

\textbf{Intui\c{c}\~ao econ\^omica (aus\^encia de distor\c{c}\~ao no topo):} o tipo $H$ recebe a cobertura eficiente ($q_H = \theta_H$) porque distorcer a cobertura do cliente de cima n\~ao ajuda a separar \textemdash{} ningu\'em ``acima'' dele tenta imit\'a-lo. J\'a o tipo $L$ \'e distorcido para baixo ($q_L = 6 < 8$) deliberadamente: encolher o pacote barato torna a imita\c{c}\~ao menos atraente para o $H$, reduzindo a renda que o operador precisa ceder. Distorce-se embaixo, preserva-se o topo.
\rubricalabel
Total 5 pontos. (i) 1{,}5 ponto por montar corretamente \emph{IR} e \emph{IC} dos dois tipos e identificar que vinculam $\textbf{(IR}_L\textbf{)}$ e $\textbf{(IC}_H\textbf{)}$ (0{,}75 por escrever as quatro restri\c{c}\~oes; 0{,}75 por apontar as duas que mordem). (ii) 2 pontos por resolver $q_H = \theta_H = 10$ (sem distor\c{c}\~ao no topo) e $q_L = \theta_L - (\theta_H - \theta_L) = 6$ (1 ponto cada). (iii) 1 ponto pelos pagamentos $T_L = 48,\ T_H = 88$ e pela renda informacional $U_H = (\theta_H-\theta_L)q_L = 12$. (iv) 0{,}5 ponto pelo coment\'ario de aus\^encia de distor\c{c}\~ao no topo (distorcer $q_H$ n\~ao relaxa nenhuma \emph{IC}, pois ningu\'em imita o tipo alto). Penalidade: distorcer $q_H$ para baixo ou $q_L$ para cima perde 1 ponto por contrariar a estrutura. Erro aritm\'etico isolado com sistema correto perde no m\'aximo 0{,}5 ponto.
\end{solucao}

\textbf{c)} (5 pontos) Compare \emph{first-best} e \emph{second-best}. \textbf{(i)} Calcule o lucro esperado do operador sob informa\c{c}\~ao assim\'etrica (\emph{second-best}) e o custo total da assimetria, isto \'e, a queda de lucro $\Pi^{FB} - \Pi^{SB}$, usando os resultados de (a) e (b). \textbf{(ii)} \emph{Sem novas contas}, nomeie e explique economicamente as \emph{duas} fontes desse custo \textemdash{} a renda informacional cedida ao tipo $H$ e a perda de efici\^encia pela distor\c{c}\~ao do tipo $L$. \textbf{(iii)} Argumente por que esse custo persistiria mesmo que o operador pudesse cobrar uma ``taxa de ades\~ao'' (\emph{lump-sum}) no in\'icio do contrato.
\resolucaobox{7.2cm}
\begin{solucao}
\textbf{Passo 1 (ponto-chave).} O custo da assimetria \'e a queda de lucro entre o \emph{first-best} e o \emph{second-best}. Sua origem \'e dupla: a renda que o operador cede ao tipo $H$ para impedir a imita\c{c}\~ao, e a perda de excedente social ao encolher o pacote do tipo $L$. Uma taxa de ades\~ao uniforme n\~ao elimina nenhuma das duas, porque o que separa os tipos \'e a \emph{(IC)}, n\~ao a \emph{(IR)}.

\textbf{Passo 2 (lucro \emph{second-best} e custo total).} Usando o menu de (b) ($q_H = 10,\ q_L = 6,\ T_H = 88,\ T_L = 48$), o lucro por tipo \'e $T_i - \tfrac{1}{2}q_i^{2}$: $88 - 50 = 38$ no tipo $H$ e $48 - 18 = 30$ no tipo $L$. Logo
\[
\Pi^{SB} = \tfrac{1}{2}\!\left[T_H - \tfrac{1}{2}q_H^{2}\right] + \tfrac{1}{2}\!\left[T_L - \tfrac{1}{2}q_L^{2}\right] = \tfrac{1}{2}\cdot 38 + \tfrac{1}{2}\cdot 30 = 34.
\]
Com $\Pi^{FB} = 41$ de (a), o custo total da assimetria \'e
\[
\Pi^{FB} - \Pi^{SB} = 41 - 34 = 7.
\]

\textbf{Passo 3 (as duas fontes do custo).} Sem exigir nova conta, o custo de $7$ tem duas origens econ\^omicas distintas. \emph{(a) Renda informacional cedida ao tipo $H$:} para impedir que o vinhedo de alta exposi\c{c}\~ao se passe por um de baixa, o operador precisa deixar-lhe excedente positivo ($U_H = (\theta_H - \theta_L)q_L = 12 > 0$); esse excedente, que no \emph{first-best} era capturado pelo operador, agora ``vaza'' para o cliente $H$. \emph{(b) Perda de efici\^encia pela distor\c{c}\~ao do tipo $L$:} ao encolher a cobertura do tipo $L$ de $q_L = 8$ (eficiente) para $q_L = 6$, o operador destr\'oi parte do excedente social do par $L$ (o valor marginal $\theta_L$ passa a exceder o custo marginal $q_L$); \'e um ganho de com\'ercio que simplesmente deixa de existir, sacrificado de prop\'osito para tornar a imita\c{c}\~ao menos atraente ao tipo $H$. A primeira fonte \'e uma \emph{transfer\^encia} (sai do operador, vai ao $H$); a segunda \'e uma \emph{perda morta} (some do sistema).

\emph{Para refer\^encia do corretor \textemdash{} n\~ao exigido do aluno:} numericamente, a renda informacional esperada \'e $\tfrac{1}{2}\,U_H = \tfrac{1}{2}\cdot 12 = 6$, e a perda de efici\^encia esperada \'e $\tfrac{1}{2}\big[(\theta_L\cdot 8 - \tfrac{1}{2}\cdot 8^{2}) - (\theta_L\cdot 6 - \tfrac{1}{2}\cdot 6^{2})\big] = \tfrac{1}{2}(32 - 30) = 1$; as duas pe\c{c}as somam $6 + 1 = 7$, exatamente o custo total.

\textbf{Passo 4 (taxa de ades\~ao).} Uma taxa fixa \emph{lump-sum} cobrada no in\'icio \'e uma transla\c{c}\~ao uniforme dos pagamentos: subtrai a mesma constante do excedente de \emph{todos} os tipos. Ela mexe na \emph{(IR)} \textemdash{} quanto excedente sobra ao cliente \textemdash{} mas n\~ao na \emph{(IC)}, que compara o excedente de \emph{ficar} no pr\'oprio pacote com o de \emph{imitar} o outro: como a mesma constante entra igualmente dos dois lados da compara\c{c}\~ao de imita\c{c}\~ao, ela cancela e a \emph{(IC)} fica intacta. Al\'em disso, $\textbf{(IR}_L\textbf{)}$ j\'a vincula ($U_L = 0$), de modo que n\~ao h\'a excedente do tipo $L$ a expropriar; se o operador elevasse a taxa para raspar a renda do $H$, violaria a $\textbf{(IR}_L\textbf{)}$ e o tipo $L$ recusaria o contrato. Como a renda do $H$ nasce da \emph{(IC)} e a distor\c{c}\~ao de $q_L$ persiste para sustent\'a-la, o custo de $7$ permanece.

\textbf{Intui\c{c}\~ao econ\^omica:} o custo da informa\c{c}\~ao privada \'e o pre\c{c}o de separar quem n\~ao se quer revelar. Parte vaza como renda ao tipo de alto valor, parte se perde ao mutilar o pacote do tipo de baixo valor para desencorajar a imita\c{c}\~ao. Ambas nascem da \emph{(IC)}, n\~ao da divis\~ao do excedente fixada pela \emph{(IR)}; por isso uma transfer\^encia sem distor\c{c}\~ao, como a taxa de ades\~ao, n\~ao as apaga.
\rubricalabel
Total 5 pontos. (i) 1{,}5 ponto por $\Pi^{SB} = 34$ (lucro esperado no \emph{second-best}, $\tfrac{1}{2}(38) + \tfrac{1}{2}(30)$). (ii) 1 ponto pelo custo total da assimetria, a queda $\Pi^{FB} - \Pi^{SB} = 41 - 34 = 7$. (iii) 1{,}5 ponto por nomear e explicar economicamente, \emph{sem exigir contas}, as duas fontes do custo: 0{,}75 pela renda informacional cedida ao tipo $H$ (transfer\^encia que o operador deixa ao $H$ para impedir a imita\c{c}\~ao) e 0{,}75 pela perda de efici\^encia da distor\c{c}\~ao do tipo $L$ (excedente social destru\'ido ao encolher $q_L$ abaixo do eficiente); n\~ao se exige a decomposi\c{c}\~ao num\'erica das parcelas. (iv) 1 ponto pelo argumento da taxa de ades\~ao: identificar que uma transfer\^encia \emph{lump-sum} entra na \emph{(IR)} mas n\~ao na \emph{(IC)} (a constante entra igualmente dos dois lados da compara\c{c}\~ao de imita\c{c}\~ao e cancela), de modo que n\~ao muda a separa\c{c}\~ao dos tipos; conceder integral apenas se ligar a persist\^encia do custo \`a \emph{(IC)} (e/ou ao fato de $\textbf{(IR}_L\textbf{)}$ j\'a vincular), meio ponto se apenas afirmar que ``n\~ao muda incentivos'' sem justificar onde entra. Erro aritm\'etico isolado com m\'etodo correto perde no m\'aximo 0{,}5 ponto (n\~ao zera o item).
\end{solucao}
